% ============================================================
%  第三卷 第六章 微扰论（§38–§45 + 习题）
%  底本：高教社中文版第6版（严肃 译，喀兴林 校）；OCR 错误已修正，脚注文本待从纸版补录
%
%  已修正的 OCR / 原书错误（均对照 原书转换/pages/p-135…p-164.png 逐条核实）：
%  1. §38 (38.7) 后正文：OCR "ε_n^(1)=0" 按原书 p123 为 c_n^(1)=0；
%  2. §38 (38.10) 后括注：OCR 之 V_mn′ 按原书 p124 为 V_mn（厄米性 V_mn=V_nm*）；
%  3. §38 (38.11)：OCR ψ_*^(0)→ψ_ν^(0)；
%  4. §38 (38.12)：第二求和式分母 OCR 作 E_n^(0)−E_k^(0)，按原书 p124 为 E_m^(0)−E_k^(0)；
%  5. §38 习题1：ω_nm 定义 OCR 之 E_n^(0)−E_n^(0) 按原书 p125 为 E_n^(0)−E_m^(0)；
%  6. §38 习题4：OCR "半轴为 ȧ=b"→a=b、x⁴ 误作 x*、"重∴Δ"→"重心"；
%     引文按原书 p125 为（А. Б. Мигдал, 1959）；ΔE_nlm 及微商两式按原书 p126–127；
%  7. §39 正文：OCR 窜入的第二处圈码①删除（该脚注文本"这个名词来自天体力学"由脚注代理转录）；
%  8. §39 习题1：OCR 丢失 ħω^(1) 记号，按原书 p129 补回（记号 ħω^(1)=√((V₁₁−V₂₂)²+4|V₁₂|²)
%     代表两个修正能级 E^(1) 之差）；
%  9. §39 习题3：Ψ₁ 式第二指数因子 OCR 漏撇号，按原书 p130 补为 e^{−(i/ħ)E^(1)′t}；
% 10. §40 (40.1)(40.2)：OCR 小写 ψ 按原书 p131 改回大写 Ψ；(40.10) OCR 漏 ±，
%     按原书 p132 为 ≠±ħω；
% 11. §40 (40.11)：OCR 之 f_kn^(0)（两处）与第一、三分母 ω_kn∓ω/ω_kn±ω 按原书 p133
%     （并经 ω→0 化回(38.12)式验证）核正为 f_km^(0) 与 ω_km∓ω、ω_km±ω；
% 12. §40 习题：(3) 式 ψ_m^(0) 按原书 p134 为 Ψ_m^(0)；"a_nΨ_n+a_nΨ_m"→a_mΨ_m；
%     概率上界 OCR "1η*1²"→|η|²；
% 13. §41：OCR 缺失的箭号补齐（t→±∞、t→−∞、t→+∞、ρ→0、λ→0）；(41.2) 指数
%     e^{iω_fi t}（OCR ω_fi′）；(41.3) 前分部积分式上限 i→t；
% 14. §41 正文："在数量级为周期的时间间隔内"按原书 p136 补 1/ω_fi；
%     "跃迁到其它态的周期"按原书为 1/ω_fi（OCR ω_n̄）；
% 15. §41：共轭方程按原书 p136 照排为 Ĥ*ψ_f*=E_fψ_f*（OCR 下标 j 系噪声）；
% 16. §41 习题1：U(x) 式末 OCR 乱码按原书 p137 为"+常数"；"分部积分 h 次"→k 次；
%     e^{ξβ/4}→e^{ξ₀²/4}；w_k0=k̄^k e^{−k̄}/k!；
% 17. §41 习题2–5 引文按原书 p138–139：（А. Б. Мигдал, 1939）、
%     （А. Б. Мигдал, 1941; Е. Л. Фейнберг, 1939）（OCR 之". .Φ"）、
%     （А. Б. Мигдал, 1941; J. S. Levinger, 1953）；"V=1/r 很小（1 与 Z 之比）"按原书 p139；
% 18. §41 习题5："只有跃迁到 l=1 的态"（OCR 重复 l=0）；dw 两式按原书 p140；
% 19. §42：OCR 乱下标 a_μ̂/a_π̂/ω_β 一律按原书 p141 为 a_fi/ω_fi；(42.4) 极限下标补 t→∞；
% 20. §43：(43.3) 指数 e^{−(i/ħ)E_i t}（OCR ε）；二级近似式中 V_fiV_ni 按原书 p144 为
%     V_fνV_νi；(43.7) V_ȷ_ν→V_fν；(43.8) V_fi′→V_fν′、(E_i′−E_ν)→(E_i−E_ν)；
% 21. §44：(44.5)(44.6) OCR 之"−"按原书 p146–147 为"~"；"三个方程式"后乱码 Φ 系噪声删除；
%     "式中的 P,E 为粒子的动量和能量，p,ε 为平面镜的动量和能量"按原书 p146（OCR 第二组误作 P）；
% 22. §45：(45.3) e^{ik′}→e^{ikr}；(45.4) OCR 乱码按原书 p148 排为"，当 ka≲1 时"；
%     (45.9) OCR 乱码按原书 p149 及推导重建为 ψ=e^{ikx}F=常数×e^{ikx}e^{−(i/ħv)∫Udx}；
% 23. §45 习题2："ψ=常数×H₀^(1)(iκr)"（OCR 乱码）按原书 p151；
% 24. 其它：OCR 粗体 \pmb{\psi}/\boldsymbol{c}/\boldsymbol{a} 一律改回斜体；"代人"→"代入"；
%     函数名收紧（\sin,\cos,\exp,\ln,\arctan）；\mathrm{d} 收紧；节内正文行首孤立的
%     "得到归一化.结果得"等断行按原书接排。
%
%  遗留问题（OCR 存疑，已照原书排版）：
%  - (43.4) 为积分路径示意图（原书以公式编号排印，MinerU 输出为图片），现用 picture
%    环境按原书 p143 重绘（E_f 实轴、极点 E_i、积分路径自极点下方绕过），编号 (43.4) 保持；
%  - §41 "(40.5)式中的积分限是这样选取的"原书如此（p135），未改；
%  - §44 "为了确定粒子的动量 P" 大写 P、§39 "我们来计算二级近似的 V_mn 矩阵元"、
%    §41 "§38 中的 c_fi"、§45 习题1 "比阱深还高一个数量级" 等表述均照原书；
%  - 脚注圈码 19 处已由脚注代理按原书页面（p137–p162）逐一转录为 \footnote{}，
%    其中 §39"久期方程"脚注"这个名词来自天体力学"按原书 p141 核录；
%  - §45 p161 脚注①中被积因子指数，原书页面印迹过小（e^{ikr} 与 e^{ik'r} 难辨），
%    按(45.3)式之 e^{ikr} 及脚注自身对 r=|x−x'| 的（无撇）定义转录为 e^{ikr}。
% ============================================================
\chapter{微扰论}
\emergencystretch 3em

\section{与时间无关的微扰}

薛定谔方程的精确解只有在少数简单情形下才能找到，对量子力学中的大多数问题来讲，所得的方程过于复杂，很难精确解出.但是往往有这样的情形，在所给问题的条件中，各量具有不同的数量级；当我们把其中较小的量略去以后，这个问题有可能变得十分简单，以致可以求出它的精确解.在这种情形下，求解这个物理问题的第一步，是求出简化问题的精确解，第二步，是计算由于在简化问题中略去了较小项而引起的误差.计算这种误差的一般方法，称为微扰论.

我们假定某一给定物理系统的哈密顿量呈以下形式：
\begin{equation*}
  \hat{H}=\hat{H}_0+\hat{V},
\end{equation*}
式中的 $\hat{V}$ 是未受扰算符 $\hat{H}_0$ 的微小修正（或称微扰项）.在 §38 和 §39 中，我们只考虑与时间无关的微扰（对 $\hat{H}_0$ 也作同样的假定）.能把 $\hat{V}$ 看作是“小于”算符 $\hat{H}_0$ 的必要条件，将在下面推导.

离散谱的微扰论问题，可按下列方式表述.假定未受扰算符 $\hat{H}_0$ 的离散谱本征值 $E_n^{(0)}$ 及其本征函数 $\psi^{(0)}$ 都是已知的，即下列方程的精确解是已知的：
\begin{equation}
  \hat{H}_0\psi^{(0)}=E^{(0)}\psi^{(0)}.
\end{equation}
现在欲求下列方程的近似解：
\begin{equation}
  \hat{H}\psi=\left(\hat{H}_0+\hat{V}\right)\psi=E\psi,
\end{equation}
也就是求受扰算符 $\hat{H}$ 的本征值 $E_n$ 和本征函数 $\psi_n$ 所具有的近似表式.

本节中，我们假定算符 $\hat{H}$ 的所有本征值都是无简并的.同时为了使结果简单起见，暂时假定本征值只有离散谱.

为方便计，最好采用矩阵形式的计算方法.为此，我们先把欲求的 $\psi$ 函数对 $\psi_n^{(0)}$ 函数组展开：
\begin{equation}
  \psi=\sum_m c_m\psi_m^{(0)},
\end{equation}
把这个展式代入(38.2)式，得
\begin{equation*}
  \sum_m c_m\left(E_m^{(0)}+\hat{V}\right)\psi_m^{(0)}=\sum_m c_mE\psi_m^{(0)},
\end{equation*}
上式两边各乘 $\psi_k^{(0)*}$ 后再积分，得
\begin{equation}
  \left(E-E_k^{(0)}\right)c_k=\sum_m V_{km}c_m.
\end{equation}
我们在这里引入了微扰算符 $\hat{V}$ 对未受扰函数组 $\psi_m^{(0)}$ 而言的矩阵 $V_{km}$：
\begin{equation}
  V_{km}=\int\psi_k^{(0)*}\hat{V}\psi_m^{(0)}\,\mathrm{d}q.
\end{equation}
我们把系数 $c_m$ 和能量 $E$ 表成以下的级数形式：
\begin{equation*}
  E=E^{(0)}+E^{(1)}+E^{(2)}+\cdots,\qquad c_m=c_m^{(0)}+c_m^{(1)}+c_m^{(2)}+\cdots,
\end{equation*}
其中的 $E^{(1)}$ 和 $c^{(1)}$ 具有与微扰 $\hat{V}$ 同样小的数量级，$E^{(2)}$ 和 $c^{(2)}$ 具有二级小的数量级，以此类推.

现在来求第 $n$ 个本征值和其本征函数的修正问题，此时，可令 $c_n^{(0)}=1$，$m\neq n$ 的 $c_m^{(0)}=0$.求第一级近似时，我们把 $E=E_n^{(0)}+E_n^{(1)}$，$c_k=c_k^{(0)}+c_k^{(1)}$ 代入(38.4)式，只保留该式中的一级小项.由 $k=n$ 的该式给出
\begin{equation}
  E_n^{(1)}=V_{nn}=\int\psi_n^{(0)*}\hat{V}\psi_n^{(0)}\,\mathrm{d}q.
\end{equation}
故本征值 $E_n^{(0)}$ 的一级修正等于微扰项在 $\psi_n^{(0)}$ 态中的平均值.

(38.4)式当 $k\neq n$ 时给出
\begin{equation}
  c_k^{(1)}=\frac{V_{kn}}{E_n^{(0)}-E_k^{(0)}},\qquad k\neq n
\end{equation}
而 $c_n^{(1)}$ 还没有确定；它必须这样选取，使得函数 $\psi_n=\psi_n^{(0)}+\psi_n^{(1)}$ 在含有一级小项为止的情形下得到归一化.根据这一点，我们须使 $c_n^{(1)}=0$.由于下列函数
\begin{equation}
  \psi_n^{(1)}=\sum_m{}'\frac{V_{mn}}{E_n^{(0)}-E_m^{(0)}}\psi_m^{(0)}
\end{equation}
（式中的撇号代表求和时除去 $m=n$ 的项）是和 $\psi_n^{(0)}$ 正交的，故 $|\psi_n^{(0)}+\psi_n^{(1)}|^2$ 的积分与 1 只差一个二级小量.

(38.8)式确定了波函数的一级近似修正.由此式可以附带地看出上述微扰法的适用条件.这个条件就是下列不等式必须成立：
\begin{equation}
  |V_{mn}|\ll|E_n^{(0)}-E_m^{(0)}|,
\end{equation}
这就是说，算符 $\hat{V}$ 的矩阵元必须远小于那两个未受扰能级间的相应间距.

其次，我们来求本征值 $E_n^{(0)}$ 在二级近似中的修正.为此，我们把 $E=E_n^{(0)}+E_n^{(1)}+E_n^{(2)}$ 和 $c_k=c_k^{(0)}+c_k^{(1)}+c_k^{(2)}$ 代入(38.4)式，观察该式中的二级小项.当 $k=n$ 时，由该式给出
\begin{equation*}
  E_n^{(2)}c_n^{(0)}=\sum_m{}'V_{nm}c_m^{(1)},
\end{equation*}
故
\begin{equation}
  E_n^{(2)}=\sum_m{}'\frac{|V_{mn}|^2}{E_n^{(0)}-E_m^{(0)}}
\end{equation}
［我们已把 $c_m^{(1)}=V_{mn}/\left(E_n^{(0)}-E_m^{(0)}\right)$ 代入，并且应用了算符 $\hat{V}$ 的厄米性：$V_{mn}=V_{nm}^{*}$.］

注意，基态能级的二级近似修正永远是负的；因为，如果 $E_n^{(0)}$ 对应于最低的能量值，则(38.10)式中所有的求和项都是负的.

更高级近似可以用同样的方法计算出来.

以上所得的结果，可以立刻推广到 $\hat{H}_0$ 算符也有连续谱的情形（但微扰仍和以前一样，应用于离散谱的一个态）.为此目的，我们只要在离散谱的求和式中加进一项连续谱的相应求积式.我们用下标 $\nu$ 来区分连续谱中的各个态，$\nu$ 的值域是连续的；我们按惯例把 $\nu$ 理解为一组物理量的数值，这组数值足以对该态进行完全描述（连续谱的态往往是简并的，这时仅仅能量值本身还不足以确定这个态），这样一来，比如(38.8)式就应该改成\footnote{这里的波函数 $\psi_\nu^{(0)}$ 应按 $\nu$ 的 $\delta$ 函数归一化.}
\begin{equation}
  \psi_n^{(1)}=\sum_m{}'\frac{V_{mn}}{E_n^{(0)}-E_m^{(0)}}\psi_m^{(0)}
  +\int\frac{V_{\nu n}}{E_n^{(0)}-E_\nu}\psi_\nu^{(0)}\,\mathrm{d}\nu.
\end{equation}
其它公式也作类似的修改.

还有一个公式值得一提，此式是用函数组 $\psi_n=\psi_n^{(0)}+\psi_n^{(1)}$ 算出物理量 $f$ 的矩阵元的微扰值，采用(38.8)式给出的 $\psi_n^{(1)}$，精确到一级项为止，不难求得如下的表式：
\begin{equation}
  f_{nm}=f_{nm}^{(0)}+\sum_k{}'\frac{V_{nk}f_{km}^{(0)}}{E_n^{(0)}-E_k^{(0)}}
  +\sum_k{}'\frac{V_{km}f_{nk}^{(0)}}{E_m^{(0)}-E_k^{(0)}}.
\end{equation}
第一个求和式中的 $k\neq n$，第二个求和式中的 $k\neq m$.

\begin{xiti}

\noindent{\heiti 习题1}\quad 求本征函数的二级近似修正 $\psi_n^{(2)}$.

解：$c_k^{(2)}$（$k\neq n$）系数可以由 $k\neq n$ 并写到二级项为止的(38.4)式算出，至于 $c_n^{(2)}$ 系数可这样选取，使 $\psi_n=\psi_n^{(0)}+\psi_n^{(1)}+\psi_n^{(2)}$ 函数在包含二级项为止的情形下得到归一化.结果得
\begin{equation*}
\begin{aligned}
  \psi_n^{(2)}={}&\sum_m{}'\sum_k{}'\frac{V_{mk}V_{kn}}{\hbar^2\omega_{nk}\omega_{nm}}\psi_m^{(0)}
  -\sum_m{}'\frac{V_{nn}V_{mn}}{\hbar^2\omega_{nm}^2}\psi_m^{(0)}-\\
  &-\frac{1}{2}\psi_n^{(0)}\sum_m{}'\frac{|V_{mn}|^2}{\hbar^2\omega_{nm}^2},
\end{aligned}
\end{equation*}
式中引进了频率
\begin{equation*}
  \omega_{nm}=\frac{1}{\hbar}\left(E_n^{(0)}-E_m^{(0)}\right)
\end{equation*}

\noindent{\heiti 习题2}\quad 求能量本征值的三级近似修正.

解：写出 $k=n$ 时(38.4)式中的三级小项，我们得
\begin{equation*}
  E_n^{(3)}=\sum_k{}'\sum_m{}'\frac{V_{nm}V_{mk}V_{kn}}{\hbar^2\omega_{mn}\omega_{kn}}
  -V_{nn}\sum_m{}'\frac{|V_{nm}|^2}{\hbar^2\omega_{mn}^2}.
\end{equation*}

\noindent{\heiti 习题3}\quad 求一非谐线性振子的能级，其哈密顿量为
\begin{equation*}
  \hat{H}=\frac{\hat{p}^2}{2m}+\frac{m}{2}\omega^2x^2+\alpha x^3+\beta x^4.
\end{equation*}

解：利用 $x$ 的矩阵元表式(23.4)，再根据矩阵的乘法规则，可以直接求出 $x^3$ 和 $x^4$ 的矩阵元.我们发现不等于零的 $x^3$ 矩阵元有
\begin{equation*}
  \left(x^3\right)_{n-3,n}=\left(x^3\right)_{n,n-3}
  =\left(\frac{\hbar}{m\omega}\right)^{3/2}\sqrt{\frac{1}{8}n(n-1)(n-2)},
\end{equation*}
\begin{equation*}
  \left(x^3\right)_{n-1,n}=\left(x^3\right)_{n,n-1}
  =\left(\frac{\hbar}{m\omega}\right)^{3/2}\sqrt{\frac{9n^3}{8}}.
\end{equation*}
这两个矩阵的对角元等于零，故哈密顿量中的 $\alpha x^3$ 项（看作谐振子的一个微扰项）在一级近似中的修正值等于零.这一项的二级近似修正和 $\beta x^4$ 的一级近似修正具有相同的数量级.$x^4$ 的对角矩阵元呈下列形式：
\begin{equation*}
  \left(x^4\right)_{n,n}=\left(\frac{\hbar}{m\omega}\right)^2\cdot\frac{3}{4}\left(2n^2+2n+1\right).
\end{equation*}
应用一般公式(38.6)和(38.10)，我们求得这个非谐振子能级的下列近似表式：
\begin{equation*}
\begin{aligned}
  E_n=\hbar\omega\left(n+\frac{1}{2}\right)
  &-\frac{15}{4}\,\frac{\alpha^2}{\hbar\omega}\left(\frac{\hbar}{m\omega}\right)^3
  \left(n^2+n+\frac{11}{30}\right)+\\
  &+\frac{3}{2}\beta\left(\frac{\hbar}{m\omega}\right)^2\left(n^2+n+\frac{1}{2}\right).
\end{aligned}
\end{equation*}

\noindent{\heiti 习题4}\quad 一个具有无限高势壁的球形势阱，经微小形变（体积不变）变成一个半轴为 $a=b$ 和 $c$ 的略长或略扁的旋转球体.求此变形阱内粒子能级的分裂（А.~Б.~Мигдал, 1959）.

解：阱壁的方程为
\begin{equation*}
  \frac{x^2+y^2}{a^2}+\frac{z^2}{c^2}=1,
\end{equation*}
通过变量变换：$x\longrightarrow ax/R$，$y\longrightarrow ay/R$，$z\longrightarrow cz/R$，上式化成 $x^2+y^2+z^2=R^2$，这是半径为 $R$ 的球方程.这个变换把粒子的哈密顿量 $\hat{H}=\hat{\boldsymbol{p}}^2/2M=-\hbar^2\Delta/2M$（$M$ 为粒子质量，能量则从阱底算起）变成 $\hat{H}=\hat{H}_0+\hat{V}$，其中
\begin{equation*}
  \hat{H}_0=-\frac{\hbar^2}{2M}\Delta,
\end{equation*}
\begin{equation*}
  \hat{V}=-\frac{\hbar^2}{2M}\left[\left(\frac{R^2}{a^2}-1\right)
  \left(\frac{\partial^2}{\partial x^2}+\frac{\partial^2}{\partial y^2}\right)
  +\left(\frac{R^2}{c^2}-1\right)\frac{\partial^2}{\partial z^2}\right].
\end{equation*}
从而把椭球阱中的运动化成了球阱中运动的问题.如果这个椭球和半径为 $R=\left(a^2c\right)^{1/3}$ 的球相差不多，则 $\hat{V}$ 可看作微扰.引入椭度 $\beta$（$|\beta|\ll1$）按下式定义
\begin{equation*}
  a\approx R\left(1-\frac{1}{3}\beta\right),\qquad c\approx R\left(1+\frac{2}{3}\beta\right),
\end{equation*}
微扰算符可写成
\begin{equation*}
  \hat{V}=\frac{\beta}{3M}\left(\hat{\boldsymbol{p}}^2-3\hat{p}_z^2\right).
\end{equation*}
一级微扰论中，粒子能级与球阱能级的差别为
\begin{equation*}
  \Delta E_{nlm}=E_{nlm}-E_{nl}^{(0)}=\langle nlm|V|nlm\rangle,
\end{equation*}
$l$ 和 $m$ 是粒子角动量及其沿椭球轴的分量；$n$ 是球阱中具有给定 $l$ 值的能级编号，它与 $m$ 无关.由于 $\boldsymbol{p}^2-3\boldsymbol{p}_z^2$ 是不可约无迹张量 $\delta_{ik}\boldsymbol{p}^2-3\boldsymbol{p}_i\boldsymbol{p}_k$ 的 $zz$ 分量，根据(107.2)和(107.6)，$\langle nlm|V|nlm\rangle$ 矩阵元正比于
\begin{equation*}
  (-1)^m\left(\begin{array}{ccc}
    l & 2 & l\\
    -m & 0 & m
  \end{array}\right),
\end{equation*}
因而
\begin{equation*}
  \langle nlm|V|nlm\rangle=\left(1-\frac{3m^2}{l(l+1)}\right)\langle nl0|V|nl0\rangle.
\end{equation*}
$3j$ 符号表见 §106.

其次，
\begin{equation*}
\begin{aligned}
  \langle nl0|V|nl0\rangle&=\frac{2}{3}\beta E_{nl}^{(0)}
  +\beta\,\frac{\hbar^2}{M}\left\langle nl0\left|\frac{\partial^2}{\partial z^2}\right|nl0\right\rangle=\\
  &=\frac{2}{3}\beta E_{nl}^{(0)}-\beta\,\frac{\hbar^2}{M}
  \int\left|\frac{\partial\psi_{nl0}}{\partial z}\right|^2r^2\,\mathrm{d}r\,\mathrm{d}o.
\end{aligned}
\end{equation*}
第一项中，我们利用了球阱中的薛定谔方程 $\hat{H}_0\psi_{nlm}=E_{nl}^{(0)}\psi_{nlm}$，第二项进行了分部积分.采用(28.11)式的 $\mathrm{Y}_{l0}$，我们求出 $\psi_{nl0}=R_{nl}(r)\mathrm{Y}_{l0}(\theta,\varphi)$ 的微商为
\begin{equation*}
  \frac{\partial}{\partial z}\psi_{nl0}
  =\left(\cos\theta\frac{\partial}{\partial r}
  -\frac{\sin\theta}{r}\frac{\partial}{\partial\theta}\right)\psi_{nl0}=
\end{equation*}
\begin{equation*}
\begin{aligned}
  ={}&-\frac{i(l+1)}{\sqrt{4(l+1)^2-1}}\left(R'_{nl}-\frac{l}{r}R_{nl}\right)\mathrm{Y}_{l+1,0}+\\
  &+\frac{il}{\sqrt{4l^2-1}}\left(R'_{nl}+\frac{l+1}{r}R_{nl}\right)\mathrm{Y}_{l-1,0}.
\end{aligned}
\end{equation*}
径部的积分用下列公式计算：
\begin{equation*}
  \int_0^{\infty}R_{nl}R'_{nl}r\,\mathrm{d}r=-\frac{1}{2}\int_0^{\infty}R_{nl}^2\,\mathrm{d}r,
\end{equation*}
\begin{equation*}
  \int_0^{\infty}R_{nl}^{\prime 2}r^2\,\mathrm{d}r=\frac{2M}{\hbar^2}E_{nl}^{(0)}
  -l(l+1)\int_0^{\infty}R_{nl}^2\,\mathrm{d}r.
\end{equation*}
这些公式是用分部积分和径向薛定谔方程(33.3)导出的：
\begin{equation*}
  R''_{nl}+\frac{2}{r}R'_{nl}-\frac{l(l+1)}{r^2}R_{nl}=-\frac{2M}{\hbar^2}E_{nl}^{(0)}.
\end{equation*}
含有 $R_{nl}^2$ 的被积项相消，最后得
\begin{equation*}
  \Delta E_{nlm}=4\beta\,\frac{l(l+1)}{(2l-1)(2l+3)}
  \left[\frac{m^2}{l(l+1)}-\frac{1}{3}\right]E_{nl}^{(0)}.
\end{equation*}
注意到
\begin{equation*}
  \frac{1}{2l+1}\sum_{m=-l}^{l}E_{nlm}=E_{nl}^{(0)},
\end{equation*}
说明多重谱线的“重心”没有位移.

\end{xiti}

\section{久期方程}

现在回到未受扰算符 $\hat{H}_0$ 的本征值具有简并时的情形.我们用 $\psi_n^{(0)},\psi_{n'}^{(0)},\cdots$ 代表属于同一能量本征值 $E_n^{(0)}$ 的一套本征函数.我们早已知道，这一套本征函数的选择方式不是唯一的；可以把它们线性组合起来，从中任意选出 $s$ 个独立的函数（$s$ 为 $E_n^{(0)}$ 能级的简并度），来代替原来的那套函数.但是，如果要求这些波函数在微扰的作用下改变都很小，那么，这套函数的选择方式就不再是任意的了.

现在我们把 $\psi_n^{(0)},\psi_{n'}^{(0)},\cdots$ 理解成为任意选定的一套未受扰本征函数.在零级近似中，正确的零级近似波函数应该是它们的线性组合，呈 $c_n^{(0)}\psi_n^{(0)}+c_{n'}^{(0)}\psi_{n'}^{(0)}+\cdots$ 形式.其中的组合系数可以和本征值的一级修正同时求出来，其步骤如下.

我们写出 $k=n,n',\cdots$ 的(38.4)式，在一级近似中，把 $E=E_n^{(0)}+E^{(1)}$ 代入；至于 $c_k$，只要用零级近似值就足够了，即 $c_n=c_n^{(0)}$，$c_{n'}=c_{n'}^{(0)}$，$\cdots$；$m\neq n,n'\cdots$ 的 $c_m=0$.于是得
\begin{equation*}
  E^{(1)}c_n^{(0)}=\sum_{n'}V_{nn'}c_{n'}^{(0)}
\end{equation*}
或
\begin{equation}
  \sum_{n'}\left(V_{nn'}-E^{(1)}\delta_{nn'}\right)c_{n'}^{(0)}=0,
\end{equation}
其中的 $n,n'$ 遍取未受扰本征值 $E_n^{(0)}$ 所属各态的各种编号值.这是一套以 $c_n^{(0)}$ 为未知数的齐次线性方程组，如果这些未知数前的系数所组成的行列式等于零，这套方程就有非零解.因此可得下列方程：
\begin{equation}
  \left|V_{nn'}-E^{(1)}\delta_{nn'}\right|=0.
\end{equation}
这是一个 $E^{(1)}$ 的 $s$ 次代数方程式，一般讲来具有 $s$ 个不同的实根.这些根就是欲求的那个本征值的各种一级修正.(39.2)式称为久期方程\footnote{这个名词来自天体力学.}.我们注意到，各根之和等于所有对角元 $V_{nn},V_{n'n'},\cdots$ 之和（这个和等于久期方程中 $\left[E^{(1)}\right]^{s-1}$ 前面的系数）.

把(39.2)式的 $s$ 个根依次代入方程组(39.1)中，解出后，可得 $s$ 套 $c_n^{(0)}$ 系数，从而求得 $s$ 个零级近似本征函数.

我们注意到，经过微扰以后，原来的简并能级一般讲来不再简并［(39.2)式一般讲来没有重根］；或者说，简并度被微扰所解除.简并度的解除可以是全部的，也可以是部分的（在后一种情况下，加入微扰后简并仍存在，但其简并度小于原来的简并度）.

也有可能发生这样的情形，由于某种原因，属于同一能级的各个简并态 $n,n',\cdots$ 之间的跃迁矩阵元全都非常小（甚至是零）.这时，我们不但要计及一级矩阵元 $V_{nn'}$，还要计及高级的不同能级之间的跃迁矩阵元 $V_{nm}\left(m\neq n,n',\cdots\right)$.我们来计算二级近似的 $V_{mn}$ 矩阵元.

在(38.4)式中 $k=n$ 时，令式左的 $E=E_n^{(0)}+E_n^{(1)}$（保留记号 $E_n^{(1)}$，作为目前近似下的能量修正值），并用 $c_n^{(0)}$ 代替 $c_n$.由于 $m\neq n,n',\cdots$ 时所有 $c_m^{(0)}=0$，我们有
\begin{equation}
  E^{(1)}c_n^{(0)}=\sum_m V_{nm}c_m^{(1)}+\sum_{n'}V_{nn'}c_{n'}^{(0)}.
\end{equation}
当 $k=m\neq n,n',\cdots$ 时，精确到一级项为止，(38.4)式，给出
\begin{equation*}
  \left[E_n^{(0)}-E_m^{(0)}\right]c_m^{(1)}=\sum_{n'}V_{mn'}c_{n'}^{(0)},
\end{equation*}
故
\begin{equation*}
  c_m^{(1)}=\sum_{n'}\frac{V_{mn'}}{E_n^{(0)}-E_m^{(0)}}c_{n'}^{(0)}.
\end{equation*}
代入(39.3)式，得
\begin{equation*}
  E^{(1)}c_n^{(0)}=\sum_{n'}c_{n'}^{(0)}
  \left(\sum_m\frac{V_{nm}V_{mn'}}{E_n^{(0)}-E_m^{(0)}}+V_{nn'}\right).
\end{equation*}
这组方程代替了(39.1)；方程组的相容条件又导致久期方程，它和(39.2)的差别在于作了下列替代：
\begin{equation}
  V_{nn'}\longrightarrow V_{nn'}+\sum_m\frac{V_{nm}V_{mn'}}{E_n^{(0)}-E_m^{(0)}}.
\end{equation}

\begin{xiti}

\noindent{\heiti 习题1}\quad 求一个两重简并能级的零级近似波函数以及本征值的一级近似修正值.

解：(39.2)式现在呈下列形式：
\begin{equation*}
  \left|\begin{array}{cc}
    V_{11}-E^{(1)} & \quad V_{12}\\
    V_{21} & \quad V_{22}-E^{(1)}
  \end{array}\right|=0
\end{equation*}
（下标 1 和 2 对应于这个两重简并能级的两个任意选定的未受扰本征函数 $\psi_1^{(0)}$ 和 $\psi_2^{(0)}$）.解之，得
\setcounter{equation}{0}\renewcommand{\theequation}{\arabic{equation}}%
\begin{equation}
  E^{(1)}=\frac{1}{2}\left[\left(V_{11}+V_{22}\right)\pm\hbar\omega^{(1)}\right],
\end{equation}
其中，记号 $\hbar\omega^{(1)}=\sqrt{\left(V_{11}-V_{22}\right)^2+4|V_{12}|^2}$，代表两个修正能级 $E^{(1)}$ 之差.再解具有以上 $E^{(1)}$ 值的(39.1)式，可得零级近似归一化波函数 $\psi^{(0)}=c_1^{(0)}\psi_1^{(0)}+c_2^{(0)}\psi_2^{(0)}$ 中的系数，其值为
\begin{equation}
\begin{aligned}
  c_1^{(0)}&=\biggl\{\frac{V_{12}}{2|V_{12}|}
  \left[1\pm\frac{V_{11}-V_{22}}{\hbar\omega^{(1)}}\right]\biggr\}^{1/2},\\
  c_2^{(0)}&=\pm\biggl\{\frac{V_{12}}{2|V_{12}|}
  \left[1\mp\frac{V_{11}-V_{22}}{\hbar\omega^{(1)}}\right]\biggr\}^{1/2}.
\end{aligned}
\end{equation}

\noindent{\heiti 习题2}\quad 推导一级近似的本征函数修正式，以及二级近似的本征值修正式.

解：我们假定 $\psi_n^{(0)}$ 已经是正确的零级近似波函数组.由这些波函数所定义的 $V_{nn'}$ 矩阵对下标 $n,n'$（属于同一简并能级的一组函数）而言显然是对角的，并且对角元 $V_{nn},V_{n'n'},\cdots$ 分别等于一级近似中的各种修正值 $E_n^{(1)},E_{n'}^{(1)},\cdots$.

现在来考虑本征函数 $\psi_n^{(0)}$ 的微扰问题.在零级近似中，$E=E_n^{(0)}$，$c_n^{(0)}=1$，$m\neq n$ 的 $c_m^{(0)}=0$.在一级近似中，我们有 $E=E_n^{(0)}+V_{nn}$，$c_n=1+c_n^{(1)}$，$c_m=c_m^{(1)}$.写出 $k\neq n,n'\cdots$ 的方程组(38.4)，保留其中的一级项：
\begin{equation*}
  \left(E_n^{(0)}-E_k^{(0)}\right)c_k^{(1)}=V_{kn}c_n^{(0)}=V_{kn},
\end{equation*}
故
\setcounter{equation}{0}%
\begin{equation}
  c_k^{(1)}=\frac{V_{kn}}{E_n^{(0)}-E_k^{(0)}},\qquad k\neq n,n',\cdots
\end{equation}
其次，我们写出 $k=n'$ 的方程式，保留其中的二级项：
\begin{equation*}
  E_n^{(1)}c_{n'}^{(1)}=V_{n'n'}c_{n'}^{(1)}+\sum_m{}'V_{n'm}c_m^{(1)}
\end{equation*}
（$m$ 的求和式中略去 $m=n,n',\cdots$ 的项）.把 $E_n^{(1)}=V_{nn}$ 以及(1)式中的 $c^{(1)}$ 代入，$n'\neq n$ 时得
\begin{equation}
  c_{n'}^{(1)}=\frac{1}{V_{nn}-V_{n'n'}}\sum_m{}'
  \frac{V_{n'm}V_{mn}}{E_n^{(0)}-E_m^{(0)}},
\end{equation}
（这个近似中的 $c_n^{(1)}$ 等于零）.\footnote{注意(1)和(2)中的各量必须很小（这是本微扰法的适用条件），这个条件只要求属于不同能级的跃迁矩阵元满足条件(38.9). 至于属于同一简并能级的跃迁矩阵元，可由久期方程精确（在一定的意义下）算出.}(1)式和(2)式确定了该组本征函数在一级近似中的修正值 $\psi_n^{(1)}=\sum_m c_m^{(1)}\psi_m^{(0)}$.

最后，写出 $k=n$ 时(38.4)式中的二级项，可得下式所示的二级能量修正：
\begin{equation}
  E_n^{(2)}=\sum_m{}'\frac{V_{nm}V_{mn}}{E_n^{(0)}-E_m^{(0)}},
\end{equation}
上式在形式上等同于(38.10)式.

\noindent{\heiti 习题3}\quad 设有一系统，在 $t=0$ 时刻处于两重简并能级的 $\psi_1^{(0)}$ 态中.在某一恒定微扰的作用下发生跃迁，求 $t$ 时刻该系统跃迁到同一能级的 $\psi_2^{(0)}$ 态的概率.

解：我们令正确的零级近似波函数为
\begin{equation*}
  \psi=c_1\psi_1+c_2\psi_2,\qquad \psi'=c_1'\psi_1+c_2'\psi_2,
\end{equation*}
其中的 $c_1,c_2$ 和 $c_1',c_2'$，就是题1的(2)式所确定的两对系数（为简单计我们把所有各量中的标号 $(0)$ 全部略去）.

反过来，有
\begin{equation*}
  \psi_1=\frac{c_2'\psi-c_2\psi'}{c_1c_2'-c_1'c_2}.
\end{equation*}
$\psi$ 和 $\psi'$ 是分属于受扰能量值 $E+E^{(1)}$ 和 $E+E^{(1)'}$ 的两个态，其中的 $E^{(1)}$ 和 $E^{(1)'}$ 是题1中的(1)式所示的两个修正值.引入时间因子后，就过渡到与时间有关的波函数：
\begin{equation*}
  \varPsi_1=\frac{\mathrm{e}^{-(i/\hbar)Et}}{c_1c_2'-c_1'c_2}
  \left[c_2'\psi\,\mathrm{e}^{-(i/\hbar)E^{(1)}t}
  -c_2\psi'\,\mathrm{e}^{-(i/\hbar)E^{(1)'}t}\right]
\end{equation*}
（$t=0$ 时，$\varPsi_1=\psi_1$）.最后，再把上式中的 $\psi$ 和 $\psi'$ 表成 $\psi_1$ 和 $\psi_2$，$\varPsi_1$ 就变成 $\psi_1$ 和 $\psi_2$ 的线性组合，其组合系数与时间有关.$\psi_2$ 函数前系数的模量平方值就是欲求的跃迁概率 $w_{21}$.把题1中的(1)式和(2)式代入，算出
\begin{equation*}
  w_{21}=\frac{2|V_{12}|^2}{\left(\hbar\omega^{(1)}\right)^2}
  \left[1-\cos\omega^{(1)}t\right].
\end{equation*}
我们看到，这个概率随时间周期性地变化，其频率为 $\omega^{(1)}$.

当时间 $t$ 远小于本题中的周期时，上式的中括号与 $t^2$ 成正比，从而 $w_{21}$ 也与 $t^2$ 成正比：
\begin{equation*}
  w_{21}=\frac{1}{\hbar^2}|V_{12}|^2t^2.
\end{equation*}
此式可用下节所给的方法［应用(40.4)］十分简单地获得.

\end{xiti}
\renewcommand{\theequation}{\arabic{section}.\arabic{equation}}%

\section{与时间有关的微扰}

现在研究明显地依赖于时间的微扰.一般讲来，我们在这种情况下谈不上本征值的修正问题，因为当哈密顿量依赖于时间时［例如受扰后的算符为 $\hat{H}=\hat{H}_0+\hat{V}(t)$］，能量不再守恒，因而也就不存在定态.我们在这里提出的问题，只是从未扰系统的定态波函数出发，对该波函数进行近似计算.

为此，我们应用一种方法，类似于求解线性微分方程组时熟知的常数变易法（P.~A.~M.~Dirac, 1926）.设 $\varPsi_k^{(0)}$（含有时间因子）为未扰系统的定态波函数.此时，未受扰波动方程的一个任意解可以表成 $\varPsi=\sum_k a_k\varPsi_k^{(0)}$ 的形式.现在来求下列受扰方程的解：
\begin{equation}
  i\hbar\frac{\partial\varPsi}{\partial t}=\left(\hat{H}_0+\hat{V}\right)\varPsi,
\end{equation}
把这个解表成以下的求和式：
\begin{equation}
  \varPsi=\sum_k a_k(t)\varPsi_k^{(0)},
\end{equation}
其中的展开系数都是时间的函数.把(40.2)式代入(40.1)式中考虑到所有的 $\varPsi_k^{(0)}$ 函数都满足下列方程：
\begin{equation*}
  i\hbar\frac{\partial\varPsi_k^{(0)}}{\partial t}=\hat{H}_0\varPsi_k^{(0)},
\end{equation*}
我们得
\begin{equation*}
  i\hbar\sum_k\varPsi_k^{(0)}\frac{\mathrm{d}a_k}{\mathrm{d}t}
  =\sum_k a_k\hat{V}\varPsi_k^{(0)}.
\end{equation*}
上式两边左乘 $\varPsi_m^{(0)*}$ 后，积分得
\begin{equation}
  i\hbar\frac{\mathrm{d}a_m}{\mathrm{d}t}=\sum_k V_{mk}(t)a_k,
\end{equation}
其中的
\begin{equation*}
\begin{aligned}
  &V_{mk}(t)=\int\varPsi_m^{(0)*}\hat{V}\varPsi_k^{(0)}\,\mathrm{d}q
  =V_{mk}\mathrm{e}^{i\omega_{mk}t},\\
  &\omega_{mk}=\frac{E_m^{(0)}-E_k^{(0)}}{\hbar}
\end{aligned}
\end{equation*}
为含有时间因子的微扰项矩阵元（必须注意，当 $\hat{V}$ 显含 $t$ 时，$V_{mk}$ 也是时间的函数）.

我们取第 $n$ 个定态的波函数作为未受扰波函数，(40.2)式中的相应系数值为 $a_n^{(0)}=1$，$k\neq n$ 的 $a_k^{(0)}=0$.求一级近似时，(40.3)式左边的 $a_k$ 写成 $a_k=a_k^{(0)}+a_k^{(1)}$，右边则用 $a_k=a_k^{(0)}$ 代入（因为其中已经包含了一级小量 $V_{mk}$）.结果得
\begin{equation}
  i\hbar\frac{\mathrm{d}a_k^{(1)}}{\mathrm{d}t}=V_{kn}(t).
\end{equation}
为了能够表明所算的修正究竟属于哪个未受扰函数，我们在 $a_k$ 系数中加进第二个下标，写成
\begin{equation*}
  \varPsi_n=\sum_k a_{kn}(t)\varPsi_k^{(0)}.
\end{equation*}
据此，我们可以把(40.4)式的积分结果写成以下的形式：
\begin{equation}
  a_{kn}^{(1)}=-\left(\frac{i}{\hbar}\right)\int V_{kn}(t)\,\mathrm{d}t
  =-\left(\frac{i}{\hbar}\right)\int V_{kn}\mathrm{e}^{i\omega_{kn}t}\,\mathrm{d}t,
\end{equation}
这个式子确定了一级近似波函数.

现在我们来更详细地考察微扰项 $\hat{V}$ 为时间的周期函数这一重要情形，$\hat{V}$ 的形式为
\begin{equation}
  \hat{V}=\hat{F}\mathrm{e}^{-i\omega t}+\hat{G}\mathrm{e}^{i\omega t},
\end{equation}
其中的 $\hat{F}$ 和 $\hat{G}$ 都是和时间无关的算符.由于 $\hat{V}$ 是厄米的，我们必须有
\begin{equation*}
  \hat{F}\mathrm{e}^{-i\omega t}+\hat{G}\mathrm{e}^{i\omega t}
  =\hat{F}^{*}\mathrm{e}^{i\omega t}+\hat{G}^{*}\mathrm{e}^{-i\omega t},
\end{equation*}
由此得 $\hat{G}=\hat{F}^{\dagger}$，即
\begin{equation}
  G_{nm}=F_{mn}^{*}.
\end{equation}
应用这个关系式，我们有
\begin{equation}
  V_{kn}(t)=V_{kn}\mathrm{e}^{i\omega_{kn}t}
  =F_{kn}\mathrm{e}^{i(\omega_{kn}-\omega)t}
  +F_{nk}^{*}\mathrm{e}^{i(\omega_{kn}+\omega)t}.
\end{equation}
代入(40.5)式并进行积分，我们得到下列波函数展开系数：
\begin{equation}
  a_{kn}^{(1)}=-\frac{F_{kn}\mathrm{e}^{i(\omega_{kn}-\omega)t}}{\hbar(\omega_{kn}-\omega)}
  -\frac{F_{nk}^{*}\mathrm{e}^{i(\omega_{kn}+\omega)t}}{\hbar(\omega_{kn}+\omega)}.
\end{equation}
这个表式只有分母都不等于零时\footnote{精确来讲，这个表式中的分母都不能很小，否则 $a_{kn}^{(1)}$ 就不能远小于 1.}才能应用，即所有的 $k$（$n$ 为给定）必须满足
\begin{equation}
  E_k^{(0)}-E_n^{(0)}\neq\pm\hbar\omega.
\end{equation}
许多应用问题中，要用到任一物理量 $f$ 对受扰波函数而言的矩阵元.一级近似中我们有
\begin{equation*}
  f_{nm}(t)=f_{nm}^{(0)}(t)+f_{nm}^{(1)}(t),
\end{equation*}
其中的
\begin{equation*}
\begin{aligned}
  f_{nm}^{(0)}(t)&=\int\varPsi_n^{(0)*}\hat{f}\varPsi_m^{(0)}\,\mathrm{d}q
  =f_{nm}^{(0)}\mathrm{e}^{i\omega_{nm}t},\\
  f_{nm}^{(1)}(t)&=\int\left[\varPsi_n^{(0)*}\hat{f}\varPsi_m^{(1)}
  +\varPsi_n^{(1)*}\hat{f}\varPsi_m^{(0)}\right]\mathrm{d}q.
\end{aligned}
\end{equation*}
把 $\varPsi_n^{(1)}=\sum_k a_{kn}^{(1)}\varPsi_k^{(0)}$ 代入上式，式中的 $a_{kn}^{(1)}$ 由(40.9)式所确定，易得所求的表式为
\begin{equation}
\begin{aligned}
  f_{nm}^{(1)}(t)=-\mathrm{e}^{i\omega_{nm}t}\sum_k\biggl\{
  &\biggl[\frac{f_{nk}^{(0)}F_{km}}{\hbar(\omega_{km}-\omega)}
  +\frac{f_{km}^{(0)}F_{nk}}{\hbar(\omega_{kn}+\omega)}\biggr]\mathrm{e}^{-i\omega t}+\\
  +&\biggl[\frac{f_{nk}^{(0)}F_{mk}^{*}}{\hbar(\omega_{km}+\omega)}
  +\frac{f_{km}^{(0)}F_{kn}^{*}}{\hbar(\omega_{kn}-\omega)}\biggr]\mathrm{e}^{i\omega t}\biggr\}.
\end{aligned}
\end{equation}
这个表式只有当每一项都不太大时才能应用，即所有的频率 $\omega_{kn},\omega_{km}$ 全都不能太接近于 $\omega$.当 $\omega=0$ 时，上式回到(38.12)式.

以上给出的各种公式，都是对只具离散谱的未受扰能级而言的.但是这些公式可以立刻推广到还具连续谱的情形中去（和以前一样，我们所考虑的微扰是对离散谱中的态而言的）；这只要在离散谱能级的求和式中，简单地补充一项对连续谱的相应求积式就可以了.这些公式中的能量 $E_k^{(0)}$，现在不但能取离散谱中的各种数值，而且能取连续谱中的各种数值，但是必须保证(40.9)(40.11)式中的分母 $\omega_{kn}\pm\omega$ 没有一个等于零.如果连续谱处于整个离散谱能级的上面，那么，在这种常遇的情况下，除了(40.10)以外，还需要补充下列条件
\begin{equation}
  E_{\min}^{(0)}-E_n^{(0)}>\hbar\omega,
\end{equation}
其中 $E_{\min}^{(0)}$ 为该连续谱的最低能级.

\begin{xiti}

\noindent{\heiti 习题}\quad 设有一周期性微扰（呈(40.6)形式），其频率 $\omega$ 满足 $E_m^{(0)}-E_n^{(0)}=\hbar(\omega+\varepsilon)$，$\varepsilon$ 是一个小量.求薛定谔方程第 $n$ 个解和第 $m$ 个解受这种周期微扰后的改变量.

解：本节所讲的方法在这里不适用，因为(40.9)式中的 $a_{mn}^{(1)}$ 系数现在变得太大.我们要从精确方程(40.3)重新出发，该式中的 $V_{mk}(t)$ 由(40.8)式给出.很明显，(40.3)式右边的各个求和项中，最重要的贡献是来自这样一些项，这些项中的时间关系由微小频率 $\omega_{mn}-\omega$ 所确定，略去所有其它各项后，我们得到以下两个联立方程：
\begin{equation*}
\begin{aligned}
  &i\hbar\frac{\mathrm{d}a_m}{\mathrm{d}t}
  =F_{mn}\mathrm{e}^{i(\omega_{mn}-\omega)t}a_n=F_{mn}\mathrm{e}^{i\varepsilon t}a_n,\\
  &i\hbar\frac{\mathrm{d}a_n}{\mathrm{d}t}=F_{mn}^{*}\mathrm{e}^{-i\varepsilon t}a_m.
\end{aligned}
\end{equation*}
我们作下列替代
\begin{equation*}
  a_n\mathrm{e}^{i\varepsilon t}=b_n,
\end{equation*}
得以下两个方程：
\begin{equation*}
  i\hbar\dot{a}_m=F_{mn}b_n,\qquad
  i\hbar\left(\dot{b}_n-i\varepsilon b_n\right)=F_{mn}^{*}a_m.
\end{equation*}
消去 $a_m$ 后，得
\begin{equation*}
  \ddot{b}_n-i\varepsilon\dot{b}_n+|F_{mn}|^2b_n/\hbar^2=0
\end{equation*}
这个方程组的两套独立解可以取成
\setcounter{equation}{0}\renewcommand{\theequation}{\arabic{equation}}%
\begin{equation}
  a_n=A\mathrm{e}^{i\alpha_1t},\qquad
  a_m=-A\frac{\hbar\alpha_1}{F_{mn}^{*}}\mathrm{e}^{i\alpha_2t}
\end{equation}
和
\begin{equation}
  a_n=B\mathrm{e}^{-i\alpha_2t},\qquad
  a_m=B\frac{\hbar\alpha_2}{F_{mn}^{*}}\mathrm{e}^{-i\alpha_1t}
\end{equation}
其中的 $A$ 和 $B$ 都是常数（它们由归一化条件确定），式中应用了下列记号：
\begin{equation*}
  \alpha_1=-\frac{\varepsilon}{2}+\Omega,\qquad
  \alpha_2=\frac{\varepsilon}{2}+\Omega,\qquad
  \Omega=\sqrt{\frac{\varepsilon^2}{4}+|\eta|^2},\qquad
  \eta=\frac{F_{mn}}{\hbar}.
\end{equation*}
因此，在所给微扰的作用下，函数 $\varPsi_n^{(0)}$，$\varPsi_m^{(0)}$ 变成函数 $a_n\varPsi_n^{(0)}+a_m\varPsi_m^{(0)}$，其中的 $a_n$ 和 $a_m$ 由(1)式及(2)式给出.

假定该系统在起始时（$t=0$）处于 $\varPsi_m^{(0)}$ 态.在随后时刻，该系统就处于我们求得的两个函数的线性组合态 $\varPsi$ 中，该组合态当 $t=0$ 时变成 $\varPsi_m^{(0)}$：
\begin{equation}
  \varPsi=\mathrm{e}^{\frac{i\varepsilon t}{2}}
  \left(\cos\Omega t-\frac{i\varepsilon}{2\Omega}\sin\Omega t\right)\varPsi_m^{(0)}
  -\frac{i\eta^{*}}{\Omega}\mathrm{e}^{-\frac{i\varepsilon t}{2}}\sin\Omega t\cdot\varPsi_n^{(0)}.
\end{equation}
$\varPsi_n^{(0)}$ 前系数的模量平方为
\begin{equation}
  \frac{|\eta|^2}{2\Omega^2}\left(1-\cos2\Omega t\right).
\end{equation}
这个式子给出了 $t$ 时刻发现该系统处于 $\varPsi_n^{(0)}$ 态的概率.我们看到这个概率是一个频率为 $2\Omega$ 的周期函数，概率值在 0 与 $|\eta|^2/\Omega^2$ 之间变化.

$\varepsilon=0$ 时（严格共振），(4)式所示的概率变成
\begin{equation*}
  \frac{1}{2}\left(1-\cos2|\eta|t\right),
\end{equation*}
它在 0 与 1 间周期性地变化；换句话说，该系统周期性地从 $\varPsi_m^{(0)}$ 态跃迁到 $\varPsi_n^{(0)}$ 态.

\end{xiti}
\renewcommand{\theequation}{\arabic{section}.\arabic{equation}}%

\section{有限时间间隔微扰作用下的跃迁}

现在假定微扰 $V(t)$ 只作用于某一有限时间间隔内［或者，当 $t\to\pm\infty$ 时，$V(t)$ 消失得足够快］.并假定该系统在微扰作用之前（或者在 $t\to-\infty$ 时）处于（离散谱的）第 $n$ 个态中.在随后时刻，系统的态将由 $\varPsi=\sum_k a_{kn}\varPsi_k^{(0)}$ 函数所确定，一级近似时，其中的系数为
\begin{equation}
\begin{aligned}
  a_{kn}&=a_{kn}^{(1)}=-\frac{i}{\hbar}\int_{-\infty}^{t}V_{kn}\mathrm{e}^{i\omega_{kn}t}\,\mathrm{d}t,
  \qquad k\neq n,\\
  a_{nn}&=1+a_{nn}^{(1)}=1-\frac{i}{\hbar}\int_{-\infty}^{t}V_{nn}\,\mathrm{d}t;
\end{aligned}
\end{equation}
(40.5)式中的积分限是这样选取的，使得 $t\to-\infty$ 时所有的 $a_{kn}^{(1)}$ 都趋于零.当微扰作用停止以后（或者在 $t\to+\infty$ 时），$a_{kn}$ 系数都取定值 $a_{kn}(\infty)$，该系统则处于下列波函数所描述的态中：
\begin{equation*}
  \varPsi=\sum_k a_{kn}(\infty)\varPsi_k^{(0)},
\end{equation*}
这个波函数又满足未受扰的波动方程，但是已不同于原来的 $\varPsi_n^{(0)}$ 函数.根据一般规则，$a_{kn}(\infty)$ 系数的模量平方值确定了该系统具有能量值 $E_k^{(0)}$ 的概率，即该系统处于第 $k$ 个定态的概率.

由此可见，在微扰作用下，该系统有可能从起初的定态过渡到任一其它定态.从初态（第 $i$ 个定态）跃迁到末态（第 $f$ 个定态）的概率等于\footnote{为一致起见，今后谈到跃迁概率时，初态和末态一律记作 $i$ 和 $f$，这个概率的下标次序为 $fi$（不是 $if$），与矩阵元的下标次序相同.}
\begin{equation}
  w_{fi}=\frac{1}{\hbar^2}\left|\int_{-\infty}^{\infty}V_{fi}\mathrm{e}^{i\omega_{fi}t}\,\mathrm{d}t\right|^2,
\end{equation}
现在来考虑一个微扰，它一经作用之后，将继续地作用下去（当然，它总是保持很小）.换句话说，$V(t)$ 当 $t\to-\infty$ 时趋于零，而当 $t\to+\infty$ 时趋向不等于零的某一有限值.这种情形下，(41.2)式不能直接应用，因为式中的积分是发散的.但是这种发散性从物理上讲来并不重要，并且很容易把它去掉.为此目的，我们进行分部积分：
\begin{equation*}
\begin{aligned}
  a_{fi}&=-\frac{i}{\hbar}\int_{-\infty}^{t}V_{fi}\mathrm{e}^{i\omega_{fi}t}\,\mathrm{d}t
  =-\left[\frac{V_{fi}\mathrm{e}^{i\omega_{fi}t}}{\hbar\omega_{fi}}\right]_{-\infty}^{t}+\\
  &+\int_{-\infty}^{t}\frac{\partial V_{fi}}{\partial t}
  \frac{\mathrm{e}^{i\omega_{fi}t}}{\hbar\omega_{fi}}\,\mathrm{d}t.
\end{aligned}
\end{equation*}
第一项的数值在下限等于零，在上限形式上与(38.8)式中的展开系数相同；式中多出一个周期因子 $\mathrm{e}^{i\omega_{fi}t}$ 的原因，仅仅是由于上式中的 $a_{fi}$ 是整个（含时）波函数 $\varPsi$ 的展开系数，而 §38 中的 $c_{fi}$ 是不含时的 $\psi$ 函数的展开系数.这就很明显，$t\to\infty$ 时这一项的极限值，不过是原函数 $\varPsi_i^{(0)}$ 在微扰的“恒定部分” $V(+\infty)$ 作用下所引起的改变，这种改变与跃迁到其它态没有什么关系.跃迁概率是由上式第二项的模量平方值给出的，并等于
\begin{equation}
  w_{fi}=\frac{1}{\hbar^2\omega_{fi}^2}
  \left|\int_{-\infty}^{\infty}\frac{\partial V_{fi}}{\partial t}
  \mathrm{e}^{i\omega_{fi}t}\,\mathrm{d}t\right|^2.
\end{equation}
对于从离散谱的态跃迁到连续谱中的态讲来，以上所得的这些公式也是成立的.其差别仅在于，当我们讲到从所给态（第 $i$ 个态）出发跃迁到值域（见 §38 末）介于 $\nu_f$ 和 $\nu_f+\mathrm{d}\nu_f$ 之间的态时，跃迁概率式(41.2)应写成下列形式：
\begin{equation}
  \mathrm{d}w_{fi}=\frac{1}{\hbar^2}
  \left|\int_{-\infty}^{\infty}V_{fi}\mathrm{e}^{i\omega_{fi}t}\,\mathrm{d}t\right|^2
  \mathrm{d}\nu_f.
\end{equation}
如果微扰 $V(t)$ 在数量级为周期 $1/\omega_{fi}$ 的时间间隔内变化得很少，则(41.2)或(41.3)式中的积分值也将很小.在极限情形下，当所加微扰变化得十分缓慢的时候，具有能量变化的（即频率 $\omega_{fi}$ 不等于零的）任何跃迁概率都趋于零.由此可见，当所加微扰变化得足够慢（或称浸渐微扰）的时候，处于某一无简并定态中的系统仍将留在该态中（还可参考 §53）.

反之，在极快地“瞬时”加入微扰这一极限情形下，“加入时刻”的导数值 $\dfrac{\partial V_{fi}}{\partial t}$ 变成无穷大.我们可以在 $\left(\dfrac{\partial V_{fi}}{\partial t}\right)\mathrm{e}^{i\omega_{fi}t}$ 的积分式中，把变化得较慢的 $\mathrm{e}^{i\omega_{fi}t}$ 因子取出积分号外，并令它等于该时刻所具有的数值.留下的积分就可以立刻算出，结果得
\begin{equation}
  w_{fi}=\frac{|V_{fi}|^2}{\hbar^2\omega_{fi}^2}.
\end{equation}
当微扰本身并不太小时，也可求出瞬时微扰下的跃迁概率.假定系统原处于哈密顿量 $\hat{H}_0$ 的某一本征函数 $\psi_i^{(0)}$ 所描述的态中.如果“瞬时”地改变哈密顿量（这就是说，所用的时间远小于原来的第 $i$ 个态跃迁到其它态的周期 $\dfrac{1}{\omega_{fi}}$），则该系统的波函数“来不及”改变，仍留于微扰以前的态中.但是，它并不是新哈密顿量 $\hat{H}$ 的本征函数，也就是说，$\psi_i^{(0)}$ 不再是定态.根据量子力学的一般规则，该系统跃迁到某一新定态的概率 $w_{fi}$ 是由 $\psi_i^{(0)}$ 函数对 $\hat{H}$ 的本征函数组 $\psi_f$ 展开后的展开系数确定的：
\begin{equation}
  w_{fi}=\left|\int\psi_i^{(0)}\psi_f^{*}\,\mathrm{d}q\right|^2.
\end{equation}
现在来证明，如果哈密顿量的改变量 $\hat{V}=\hat{H}-\hat{H}_0$ 很小，这个一般公式如何化回到(41.5)式.以下两式
\begin{equation*}
  \hat{H}_0\psi_i^{(0)}=E_i^{(0)}\psi_i^{(0)},\qquad
  \hat{H}^{*}\psi_f^{*}=E_f\psi_f^{*}
\end{equation*}
分别各乘上 $\psi_f^{*}$ 和 $\psi_i^{(0)}$ 后，对 $\mathrm{d}q$ 积分，并逐项相减；再利用 $\hat{H}$ 算符的厄米性，可得
\begin{equation*}
  \left(E_f-E_i^{(0)}\right)\int\psi_f^{*}\psi_i^{(0)}\,\mathrm{d}q
  =\int\psi_f^{*}\hat{V}\psi_i^{(0)}\,\mathrm{d}q.
\end{equation*}
如果微扰 $\hat{V}$ 很小，则一级近似下的 $E_f$ 可用未微扰能级 $E_f^{(0)}$ 代替，而波函数 $\psi_f$（在上式右边的）相应地可用 $\psi_f^{(0)}$ 代替.此时得
\begin{equation*}
  \int\psi_f^{*}\psi_i^{(0)}\,\mathrm{d}q
  =\frac{1}{\hbar\omega_{fi}}\int\psi_f^{(0)*}\hat{V}\psi_i^{(0)}\,\mathrm{d}q.
\end{equation*}
故(41.6)式变成了(41.5)式.

\begin{xiti}

\noindent{\heiti 习题1}\quad 处于基态的一个带电振子上突然加入一个均匀电场，求该振子受此微扰后跃迁到激发态的概率.

解：均匀电场中的振子势能（电场的作用力为 $F$）为
\begin{equation*}
  U(x)=\frac{m\omega^2}{2}x^2-Fx
  =\frac{m\omega^2}{2}\left(x-x_0\right)^2+\text{常数}
\end{equation*}
（其中的 $x_0=\dfrac{F}{m\omega^2}$），故仍具振子形式（但平衡位置移动了）.因此受扰振子的定态波函数为 $\psi_k\left(x-x_0\right)$，而 $\psi_k(x)$ 是(23.12)式的振子波函数；初态波函数是(23.13)式的 $\psi_0(x)$.根据这些公式以及(23.11)式的厄米多项式，可得
\begin{equation*}
  \int_{-\infty}^{\infty}\psi_0^{(0)}\psi_k\,\mathrm{d}x
  =\frac{(-1)^k}{\sqrt{2^k\pi k!}}\,\mathrm{e}^{-\xi_0^2/2}
  \int_{-\infty}^{\infty}\mathrm{e}^{-\xi\xi_0}
  \frac{\mathrm{d}^k}{\mathrm{d}\xi^k}\mathrm{e}^{-\xi^2+2\xi\xi_0}\,\mathrm{d}\xi,
\end{equation*}
式中引入了记号 $\xi_0=x_0\sqrt{\dfrac{m\omega}{\hbar}}$.分部积分 $k$ 次以后，上式右边积分变成下列形式：
\begin{equation*}
  \xi_0^k\int_{-\infty}^{\infty}\mathrm{e}^{-\xi^2+\xi\xi_0}\,\mathrm{d}\xi
  =\xi_0^k\sqrt{\pi}\,\mathrm{e}^{\xi_0^2/4}.
\end{equation*}
故跃迁概率(41.6)为
\begin{equation*}
  w_{k0}=\frac{\bar{k}^k}{k!}\,\mathrm{e}^{-\bar{k}},\qquad
  \bar{k}=\frac{1}{2}\xi_0^2=F^2/2m\hbar\omega^3.
\end{equation*}
作为量子数 $k$ 的函数，上式呈泊松分布形式，其中 $\bar{k}$ 是 $k$ 的平均值.

当 $F$ 很小，以致 $\bar{k}\ll1$ 时，微扰论可以适用.此时，跃迁到激发态的概率很小，并随 $k$ 的增大而很快衰减.最大值为 $w_{10}\approx\bar{k}$.

相反的情形下，当 $F$ 很大时（$\bar{k}\gg1$），激发振子的产生具有极大的概率；该振子仍留于基态的概率为 $w_{00}=\mathrm{e}^{-\bar{k}}$.

\noindent{\heiti 习题2}\quad 处于基态的原子核，受一急剧撞击后具有速度 $v$；撞击时间 $\tau$ 假定既小于电子周期，又小于 $a/v$（$a$ 为原子线度）.求该原子在这种“冲击”影响下的激发概率（А.~Б.~Мигдал, 1939）.

解：变换到 $K'$ 参考系，这个参考系与受撞后的原子核一起运动.由于条件 $\tau\ll a/v$，受撞期间的原子核可以认为实际上没有移动，因此，受扰刚结束时，参考系 $K'$ 中的电子坐标与原参考系 $K$ 中的电子坐标相等.初态波函数在 $K'$ 参考系中为
\begin{equation*}
  \psi_0'=\psi_0\exp\left(-i\boldsymbol{q}\cdot\sum_a\boldsymbol{r}_a\right),
  \qquad\boldsymbol{q}=\frac{m\boldsymbol{v}}{\hbar},
\end{equation*}
其中的 $\psi_0$ 为原子核未动时的基态波函数，指数幂中的 $\sum_a$ 是对该原子中所有的 $Z$ 个电子求和.按(41.6)式，跃迁到第 $k$ 个激发态的概率由下式确定：
\begin{equation*}
  w_{k0}=\left|\left\langle k\left|
  \exp\left(-i\boldsymbol{q}\cdot\sum_a\boldsymbol{r}_a\right)
  \right|0\right\rangle\right|^2.
\end{equation*}
特别是当 $qa\ll1$ 时，积分号内的指数因子可以展开.并由于 $\psi_0$ 和 $\psi_k$ 函数的正交性，注意 $\psi_k^{*}\psi_0$ 的积分等于零，故得
\begin{equation*}
  w_{k0}=\left|\left\langle k\left|
  \left(\boldsymbol{q}\cdot\sum_a\boldsymbol{r}_a\right)
  \right|0\right\rangle\right|^2.
\end{equation*}

\noindent{\heiti 习题3}\quad 求氢原子受突然“冲击”后激发和电离的总概率（参考上题）.

解：欲求的总概率可按下式计算：
\begin{equation*}
  1-w_{00}=1-\left|\int\psi_0^2\mathrm{e}^{-i\boldsymbol{q}\cdot\boldsymbol{r}}\,\mathrm{d}V\right|^2,
\end{equation*}
其中的 $w_{00}$ 为原子仍留于基态的概率（氢原子基态波函数
$\psi_0=\dfrac{1}{\sqrt{\pi a^3}}\,\mathrm{e}^{-r/a}$；$a$ 为玻尔半径）.算出积分后得
\begin{equation*}
  1-w_{00}=1-\frac{1}{\left(1+\dfrac{1}{4}q^2a^2\right)^4}.
\end{equation*}
在 $qa\ll1$ 的极限情形下，这个概率变成 $1-w_{00}\approx q^2a^2$ 而趋于零，当 $qa\gg1$ 时，$1-w_{00}\approx1-\left(2/qa\right)^8$ 趋于 1.

\noindent{\heiti 习题4}\quad 原子序数 $Z$ 很大的原子核进行 $\beta$ 衰变，求该原子 $K$ 层电子的逸出概率.$\beta$ 粒子的速度假定大于 $K$ 电子的速度（А.~Б.~Мигдал, 1941; Е.~Л.~Фейнберг, 1939）.

解\footnote{题 4 和题 5 中应用了原子单位.}：在上述条件下，$\beta^-$ 粒子穿过 $K$ 层的时间小于电子的绕核周期，因而核电荷的变化可以看作是瞬时的.电荷变化所引起的核电场改变量 $V=1/r$ 很小（1 与 $Z$ 之比），起着微扰的作用.按(41.5)式，能量\footnote{对 $K$ 电子，今后将用类氢原子态（参考 §74）.}为 $E_0=-Z^2/2$ 的两个 $K$ 层电子中，有一个电子跃迁到能量介于 $E=k^2/2$ 和 $E+\mathrm{d}E$（$\mathrm{d}E=k\,\mathrm{d}k$）的连续谱区间内的概率为
\begin{equation*}
  \mathrm{d}w=2\,\frac{4|V_{0k}|^2}{\left(k^2+Z^2\right)^2}\,\mathrm{d}k.
\end{equation*}
在决定矩阵元 $V_{0k}$ 的表式时，近核处（$\sim1/Z$）的积分区间最重要，这时连续谱状态波函数仍可利用类氢原子的表式.末态电子的角动量必为 $l=0$（等于初态的角动量）.利用 §36 得到的 $R_{10}$ 和按 $k/2\pi$ 标度归一化的 $R_{k0}$ 函数以及数学附录中的(f.3)式，得\footnote{计算时最好采用库仑单位. 最后结果中再换成原子单位.}
\begin{equation*}
  \left(\frac{1}{r}\right)_{0k}
  =\frac{4\sqrt{2\pi k}\left(1+\dfrac{ik}{Z}\right)^{iZ/k}
  \left(1-\dfrac{ik}{Z}\right)^{-iZ/k}}
  {\left(1-\mathrm{e}^{-2\pi Z/k}\right)\left(1+\dfrac{k^2}{Z^2}\right)}.
\end{equation*}
由于
\begin{equation*}
  \left|\left(1+i\alpha\right)^{i/\alpha}\right|^2
  =\exp\left(-2\,\frac{\arctan\alpha}{\alpha}\right),
\end{equation*}
最后得
\begin{equation*}
  \mathrm{d}w=\frac{2^7}{Z^4\left(1+\dfrac{k^2}{Z^2}\right)^4}
  f\left(\frac{k}{Z}\right)k\,\mathrm{d}k,
\end{equation*}
并有
\begin{equation*}
  f(\alpha)=\frac{1}{1-\mathrm{e}^{-2\pi/\alpha}}
  \exp\left(-4\,\frac{\arctan\alpha}{\alpha}\right).
\end{equation*}
$f(\alpha)$ 函数的极限值为
\begin{equation*}
  \alpha\ll1\ \text{时}\ f=\mathrm{e}^{-4},\qquad
  \alpha\gg1\ \text{时}\ f=\alpha/2\pi.
\end{equation*}
$\mathrm{d}w$ 对逸出电子的所有能量积分，可得 $K$ 层电离的总概率.数值计算后，得 $w=0.65Z^2$.

\noindent{\heiti 习题5}\quad 原子序数 $Z$ 很大的原子核进行 $\alpha$ 衰变，求该原子 $K$ 层电子的逸出概率.$\alpha$ 粒子的速度小于 $K$ 电子的速度，但是它从原子核内逸出的时间小于电子的绕核周期（А.~Б.~Мигдал, 1941; J.~S.~Levinger, 1953）.

解：$\alpha$ 粒子刚逸出时，作用于电子的微扰还是浸渐微扰.因此，所求效应主要取决于浸渐被破坏那一微扰“加入时刻”附近的时间间隔，此时的 $\alpha$ 粒子虽已逸出原子核并作自由运动，但仍处于比 $K$ 层电子的轨道半径来得小的距离内.$\alpha$ 粒子和原子核的联合场与纯库仑场 $Z/r$ 的差别，代表了致使原子电离的微扰作用 $V$.原子量分别为 4 和 $A-4$，电荷分别为 2 和 $Z-2$ 相互距离为 $vt$（$v$ 是 $\alpha$ 粒子和原子核的相对速度）的两个粒子的偶极矩等于
\begin{equation*}
  \frac{2\left(A-4\right)-4\left(Z-2\right)}{A}vt=\frac{2\left(A-2Z\right)}{A}vt.
\end{equation*}
故原子核和 $\alpha$ 粒子的势场的偶极项为\footnote{如果 $A-2Z$ 很小，还需考虑下一项四极矩项.}
\begin{equation*}
  V=\frac{2\left(A-2Z\right)}{A}vt\,\frac{z}{r^3},
\end{equation*}
式中取 $z$ 轴沿速度 $\boldsymbol{v}$ 的方向.取电子运动方程 $\ddot{z}=-Zz/r^3$ 的矩阵元，可把以上微扰 $V$ 的矩阵元化成 $z$ 的矩阵元，我们有
\begin{equation*}
  \left(\frac{z}{r^3}\right)_{0k}=\frac{\left(E-E_0\right)^2}{Z}z_{0k}.
\end{equation*}
按(41.2)式，两个 $K$ 电子中有一个电子进行跃迁的概率为
\begin{equation*}
  \mathrm{d}w=2\left|\int_0^{\infty}V_{0k}\mathrm{e}^{i(E_0-E)t}\,\mathrm{d}t\right|^2
  \mathrm{d}k=\frac{8\left(A-2Z\right)^2v^2}{A^2Z^2}|z_{0k}|^2\,\mathrm{d}k/2\pi.
\end{equation*}
计算此积分时，先在被积函数中引进一个阻尼因子 $\mathrm{e}^{-\lambda t}$，$\lambda>0$，积出后，再令 $\lambda\to0$，计算 $z=r\cos\theta$ 的矩阵元时，由于初态的轨道角动量 $l=0$，因此只有跃迁到 $l=1$ 的态时 $\cos\theta$ 的矩阵元才不等于零；故有
\begin{equation*}
  \left|\left(\cos\theta\right)_{01}\right|^2
  =\left(\cos^2\theta\right)_{00}=\frac{1}{3},\qquad
  |z_{0k}|^2=\frac{1}{3}|r_{0k}|^2.
\end{equation*}
应用径函数 $R_{00}$ 和 $R_{k1}$ 算出 $r_{0k}$ 后，结果得
\begin{equation*}
  \mathrm{d}w=\frac{2^{11}\left(A-2Z\right)^2v^2}
  {3A^2Z^6\left(1+\dfrac{k^2}{Z^2}\right)^5}
  f\left(\frac{k}{Z}\right)k\,\mathrm{d}k.
\end{equation*}
$f$ 函数已见题4.

\end{xiti}

\section{周期微扰作用下的跃迁}

在周期微扰的作用下跃迁到连续谱状态的概率，和以前的结果有所不同.设在 $t=0$ 的起初时刻，系统处于离散谱的第 $i$ 个定态中.假定周期微扰的频率 $\omega$ 满足
\begin{equation}
  \hbar\omega>E_{\min}-E_i^{(0)},
\end{equation}
其中的 $E_{\min}$ 为连续谱的起始能量值.

根据 §40 的结果显然可知，起首要作用的态是能量 $E_f$ 十分接近于共振能量 $E_i^{(0)}+\hbar\omega$ 的那些态，亦即 $\omega_{fi}-\omega$ 值很小的那些态.根据这一点，(40.8)式的微扰矩阵元中只要考虑（频率 $\omega_{fi}-\omega$ 接近于零的）第一项就足够了.把这一项代入(40.5)式中进行积分，我们得
\begin{equation}
  a_{fi}=-\frac{i}{\hbar}\int_0^{t}V_{fi}(t)\,\mathrm{d}t
  =-F_{fi}\frac{\mathrm{e}^{i(\omega_{fi}-\omega)t}-1}{\hbar(\omega_{fi}-\omega)}.
\end{equation}
积分下限的选定，要使 $t=0$ 时 $a_{fi}=0$，以便符合于所设的初始条件.

由此求得的 $a_{fi}$ 模量平方值为
\begin{equation}
  |a_{fi}|^2=|F_{fi}|^2\cdot
  \frac{4\sin^2\dfrac{\omega_{fi}-\omega}{2}\,t}{\hbar^2\left(\omega_{fi}-\omega\right)^2}.
\end{equation}
很易看出，当 $t$ 很大时，上式可以看作和 $t$ 成正比的一个函数.证明时，只要注意到
\begin{equation}
  \lim_{t\to\infty}\frac{\sin^2\alpha t}{\pi t\alpha^2}=\delta(\alpha).
\end{equation}
实际上，$\alpha\neq0$ 时这个极限值等于零，但当 $\alpha=0$ 时我们有
$\dfrac{\sin^2\alpha t}{t\alpha^2}=t$，故其极限值为无穷大；最后，对 $\alpha$ 从 $-\infty$ 到 $+\infty$ 积分，我们有（作替代 $\alpha t=\xi$）
\begin{equation*}
  \frac{1}{\pi}\int_{-\infty}^{\infty}\frac{\sin^2\alpha t}{t\alpha^2}\,\mathrm{d}\alpha
  =\frac{1}{\pi}\int_{-\infty}^{\infty}\frac{\sin^2\xi}{\xi^2}\,\mathrm{d}\xi=1.
\end{equation*}
可见(42.4)式左边的函数实际上满足 $\delta$ 函数定义中的所有条件.故 $t$ 很大时我们有
\begin{equation*}
  |a_{fi}|^2=\frac{1}{\hbar^2}|F_{fi}|^2\pi t\,
  \delta\left(\frac{\omega_{fi}-\omega}{2}\right),
\end{equation*}
以 $\hbar\omega_{fi}=E_f-E_i^{(0)}$ 代入，并应用 $\delta(ax)=\delta(x)/a$，得
\begin{equation*}
  |a_{fi}|^2=\frac{2\pi}{\hbar}|F_{fi}|^2
  \delta\left(E_f-E_i^{(0)}-\hbar\omega\right)t.
\end{equation*}
表式 $|a_{fi}|^2\mathrm{d}\nu_f$ 等于从原态跃迁到 $\mathrm{d}\nu_f$ 间隔内的一个态的概率.我们看到，当 $t$ 很大时，它与从 $t=0$ 开始所经历的时间间隔成正比.单位时间内的跃迁概率 $\mathrm{d}w_{fi}$ 为\footnote{很易证明，当我们考虑了(40.8)式中曾被略去的第二项以后，得到一个补充项，这一项除以 $t$，当 $t\to\infty$ 时趋于零.}
\begin{equation}
  \mathrm{d}w_{fi}=\frac{2\pi}{\hbar}|F_{fi}|^2
  \delta\left(E_f-E_i^{(0)}-\hbar\omega\right)\mathrm{d}\nu_f.
\end{equation}
如所预料，除了跃迁到能量为 $E_f=E_i^{(0)}+\hbar\omega$ 的态外，这个概率等于零.如果该连续谱不存在简并能级，则 $\nu_f$ 可以径自取作能量值，此时 $\mathrm{d}\nu_f$ “间隔”内的态就变成能量为 $E=E_i^{(0)}+\hbar\omega$ 的一个单态，跃迁到这个态的概率为
\begin{equation}
  w_{Ei}=\frac{2\pi}{\hbar}|F_{Ei}|^2.
\end{equation}
还有另一种方法能导出(42.5)式，它在方法论上很有教益，此法中并不假定周期微扰是在 $t=0$ 时刻加入的，而是从 $t=-\infty$ 开始以指数规律 $\mathrm{e}^{\lambda t}$ 逐渐地加入（浸渐加入），$\lambda$ 是一个正的常数，以后让它趋于零.初始条件 $a_{fi}=0$ 现在适用于 $t=-\infty$ 时刻，微扰矩阵元变成
\begin{equation*}
  V_{fi}(t)=F_{fi}\,\mathrm{e}^{i(\omega_{fi}-\omega)t+\lambda t},
\end{equation*}
(42.2)式变成
\begin{equation}
  a_{fi}=-\frac{i}{\hbar}\int_{-\infty}^{t}V_{fi}(t)\,\mathrm{d}t
  =-F_{fi}\frac{\mathrm{e}^{i(\omega_{fi}-\omega)t+\lambda t}}
  {\hbar\left(\omega_{fi}-\omega-i\lambda\right)}.
\end{equation}
因而
\begin{equation*}
  |a_{fi}|^2=\frac{1}{\hbar^2}|F_{fi}|^2
  \frac{\mathrm{e}^{2\lambda t}}{\left(\omega_{fi}-\omega\right)^2+\lambda^2},
\end{equation*}
每单位时间的跃迁概率由上式微商给出：
\begin{equation*}
  \frac{\mathrm{d}}{\mathrm{d}t}|a_{fi}|^2=2\lambda|a_{fi}|^2,
\end{equation*}
和(42.4)式一样，下式是成立的：
\begin{equation}
  \lim_{\lambda\to0}\frac{\lambda}{\pi\left(\alpha^2+\lambda^2\right)}=\delta(\alpha),
\end{equation}
利用上式取极限 $\lambda\to0$，可得
\begin{equation*}
  \frac{\mathrm{d}}{\mathrm{d}t}|a_{fi}|^2\longrightarrow
  \frac{2\pi}{\hbar^2}|F_{fi}|^2\delta(\omega_{fi}-\omega),
\end{equation*}
这就回到了(42.5)式.

\section{连续谱中的跃迁}

微扰论的一个最重要应用，是计算恒定（与时间无关）微扰作用下连续谱中的跃迁概率.我们早就指出过，连续谱的态差不多总是简并的.用某种方式选好对应于某一给定能级的一套未受扰波函数以后，我们可以把这个问题作如下处理：设在某一初始时刻，已知该系统处于其中的某一态中，试求跃迁到具有同一能量的另一态的概率.对从初态 $i$ 跃迁到 $\nu_f$ 和 $\nu_f+\mathrm{d}\nu_f$ 间隔内的态讲来，根据(42.5)式我们立刻得（改变式中的记号，并令该式中的 $\omega=0$）
\begin{equation}
  \mathrm{d}w_{fi}=\frac{2\pi}{\hbar}|V_{fi}|^2\delta(E_f-E_i)\,\mathrm{d}\nu_f.
\end{equation}
这个表式正如我们所预料的，除了 $E_f=E_i$ 以外都等于零：在恒定微扰的作用下，只有能量相同的态才能彼此跃迁.必须注意的是，对连续谱状态的跃迁而言，$\mathrm{d}w_{fi}$ 不能直接看作跃迁概率；甚至它的量纲也不一定是［1/时间］.在(43.1)式中，$\mathrm{d}w_{fi}$ 代表单位时间内的跃迁次数，其量纲依赖于连续谱波函数所选的归一化方式.\footnote{此处所讲的理论包括很多现象，例如各种碰撞；系统无论在初态和末态都是一群自由粒子，而微扰就是它们之间的相互作用. 波函数适当归一化后，(43.1)式就成为碰撞截面（见 §126）.}

我们来计算受扰后的波函数，它在微扰作用之前等于原来的未受扰波函数 $\psi_i^{(0)}$.采用 §42 末段所给的方法，我们可把微扰看作是按 $\mathrm{e}^{\lambda t}$ 规律浸渐地加入的，然后取 $\lambda\to0$.按(42.7)式，令 $\omega=0$，并更换记号，得
\begin{equation}
  a_{fi}^{(1)}=V_{fi}\frac{\exp\left\{\dfrac{i}{\hbar}(E_f-E_i)t+\lambda t\right\}}
  {E_i-E_f+i\lambda\hbar}.
\end{equation}
受扰波函数为
\begin{equation*}
  \varPsi_i=\varPsi_i^{(0)}+\int a_{fi}^{(1)}\varPsi_f^{(0)}\,\mathrm{d}\nu_f.
\end{equation*}
式中的积分遍及整个连续谱\footnote{如果还有离散谱，则在此式（以及随后的式子）中尚需加进对离散谱各态的求和项.}.把(43.2)式代入，得
\begin{equation}
  \varPsi_i=\left[\psi_i^{(0)}+\int V_{fi}\psi_f^{(0)}
  \frac{\mathrm{d}\nu_f}{E_i-E_f+i0}\right]
  \exp\left(-\frac{i}{\hbar}E_it\right).
\end{equation}
取极限 $\lambda\to0$ 时，$\mathrm{e}^{\lambda t}$ 因子变成 1.$+i0$ 项代表 $\lambda$ 从正值趋于零时 $i\lambda$ 的极值，它确定了对变量 $E_f$ 的积分方式（$\mathrm{d}E_f$ 是 $\mathrm{d}\nu_f$ 中的一个因子，它和描述连续谱态的其它量的微分乘在一起）.没有 $i\lambda$ 项，(43.3)的被积函数在 $E_f=E_i$ 处就会有极点，积分在极点附近将会发散，$i\lambda$ 项把这个极点移到复变量 $E_f$ 的上半平面中去，作完 $\lambda\to0$ 的极限手续后，极点回到实轴上，但是我们知道，此时积分路径必须从极点的下面绕过去：
\begin{equation}
  \setlength{\unitlength}{0.9pt}%
  \begin{picture}(152,34)(0,0)
    \put(2,16){\line(1,0){24}}
    \put(28,16){\vector(1,0){14}}
    \put(44,16){\line(1,0){18}}
    \qbezier(62,16)(62,5)(73,5)
    \qbezier(73,5)(84,5)(84,16)
    \put(86,16){\vector(1,0){16}}
    \put(104,16){\line(1,0){44}}
    \put(73,16){\circle*{2}}
    \put(73,16){\makebox(0,0)[b]{$E_i$}}
    \put(148,16){\makebox(0,0)[t]{$E_f$}}
  \end{picture}
\end{equation}
(43.3)式中的时间因子表明，这个函数和原来的未受扰函数属于同一能量 $E_i$.换句话说，下列函数
\begin{equation}
  \psi_i=\psi_i^{(0)}+\int\frac{V_{fi}}{E_i-E_f+i0}\psi_f^{(0)}\,\mathrm{d}\nu_f
\end{equation}
满足薛定谔方程
\begin{equation*}
  \left(\hat{H}_0+\hat{V}\right)\psi_i=E_i\psi_i.
\end{equation*}
因此，所得表式自然和(38.8)式一致\footnote{采用此式时，根据 $\psi_i$ 在远距离处的渐近式中只能包含出射波（不能含入射波）这一条件，就能确定积分路径（见 §136）.}.

上述计算相当于微扰论的一级近似.不难算出它的二级近似，为此，我们需要推导 $\varPsi_i$ 的下一级近似表式.采用 §38 的方法（现在我们已经知道了处理“发散”积分的方法）容易做到这一点.简单计算后，得
\begin{equation*}
  \varPsi_i=\biggl\{\psi_i^{(0)}+\int\biggl[V_{fi}
  +\int\frac{V_{f\nu}V_{\nu i}}{E_i-E_\nu+i0}\,\mathrm{d}\nu\biggr]\times
  \frac{\psi_f^{(0)}\,\mathrm{d}\nu_f}{E_i-E_f+i0}\biggr\}
  \mathrm{e}^{-(i/\hbar)E_it}.
\end{equation*}
将此式和(43.3)式比较，直接模拟(43.1)式，即能写出相应的跃迁概率（确切地说，是跃迁次数）公式：
\begin{equation}
  \mathrm{d}w_{fi}=\frac{2\pi}{\hbar}
  \left|V_{fi}+\int\frac{V_{f\nu}V_{\nu i}}{E_i-E_\nu+i0}\,\mathrm{d}\nu\right|^2
  \delta(E_i-E_f)\,\mathrm{d}\nu_f.
\end{equation}
矩阵元 $V_{fi}$ 对于所考虑的跃迁讲来，有可能等于零.此时一级近似为零，(43.6)式变成
\begin{equation}
  \mathrm{d}w_{fi}=\frac{2\pi}{\hbar}
  \left|\int\frac{V_{f\nu}V_{\nu i}}{E_i-E_\nu}\,\mathrm{d}\nu\right|^2
  \delta(E_i-E_f)\,\mathrm{d}\nu_f.
\end{equation}
应用上式时，$E_\nu=E_i$ 点一般不是被积函数的极点，于是对 $E_\nu$ 的积分方式就无关紧要了，积分可沿实轴进行.

$V_{f\nu}$ 和 $V_{\nu i}$ 都不等于零的那些 $\nu$ 态，通常称作 $i\to f$ 跃迁的中间态.直观地讲，好像这种跃迁分成 $i\to\nu$ 和 $\nu\to f$ 两步走（当然，这只是一种口头上的说法）.也有可能 $i\to f$ 的跃迁不能通过一个而要连接通过好几个中间态才能发生.(43.7)式可以立即推广到这种情形，例如，如果需要两个中间态，则有
\begin{equation}
  \mathrm{d}w_{fi}=\frac{2\pi}{\hbar}
  \left|\int\frac{V_{f\nu'}V_{\nu'\nu}V_{\nu i}}
  {\left(E_i-E_{\nu'}\right)\left(E_i-E_\nu\right)}
  \,\mathrm{d}\nu\,\mathrm{d}\nu'\right|^2
  \delta(E_f-E_i)\,\mathrm{d}\nu_f.
\end{equation}
最后，为了澄清取(43.4)的积分路径的数学含义，我们来证明下列公式：
\begin{equation}
  \int\frac{f(x)\,\mathrm{d}x}{x-a-i0}
  =P\int\frac{f(x)\,\mathrm{d}x}{x-a}+i\pi f(a),
\end{equation}
式中积分沿着包括 $x=a$ 点在内的实轴段.如果我们沿着半径为 $\rho$ 的一个半圆绕过极点 $x=a$，我们发现，整个积分等于沿实轴从下限积到 $x=a-\rho$ 再从 $a+\rho$ 积到上限的两个积分，再加上被积函数在极点处的留数的 $i\pi$ 倍.取极限 $\rho\to0$ 时，沿实轴的两个积分就变成沿整段的主值积分（用 $P$ 表示），结果就是(43.9)式，该式也可用符号形式表成
\begin{equation}
  \frac{1}{x-a-i0}=P\,\frac{1}{x-a}+i\pi\delta(x-a),
\end{equation}
$P$ 代表对 $f(x)/(x-a)$ 的积分取主值.

\section{能量的不确定度关系}

我们来考虑由相互作用很弱的两个部分所组成的一个系统.假定在某一时刻，已知这两个部分都具有确定的能量值，我们把它记作 $E$ 和 $\varepsilon$.经过某一时间间隔 $\Delta t$ 以后，再测能量；所得的 $E',\varepsilon'$ 值一般讲来和 $E,\varepsilon$ 并不相同.容易求出测量结果之差 $E'+\varepsilon'-E-\varepsilon$ 的最概然值所具有的数量级.

根据 $\omega=0$ 的(42.3)式，系统在不含时的微扰作用下，经过时间 $t$ 后从能量为 $E$ 的态跃迁到能量为 $E'$ 的态所具有的跃迁概率与下式成正比：
\begin{equation*}
  \frac{\sin^2\left(\dfrac{E'-E}{2\hbar}\,t\right)}{\left(E'-E\right)^2}
\end{equation*}
由此可知，差 $E'-E$ 的最概然值具有 $\hbar/t$ 的数量级.

把这个结果用到目前考虑的情形中（以两个子系统间的相互作用为微扰），我们得到以下的关系式：
\begin{equation}
  \left|E+\varepsilon-E'-\varepsilon'\right|\Delta t\sim\hbar.
\end{equation}
由此可见，时间间隔 $\Delta t$ 愈短，观测到的能量变化就愈大.重要的是，$\hbar/\Delta t$ 这个数量级和微扰本身的数值无关.不管子系统间的相互作用弱到什么地步，由(44.1)式确定的能量变化值都能观察到.这个结果是量子理论所特有的，并且具有深刻的物理意义.它表明，通过两次测量来验证量子力学中的能量守恒定律只有 $\hbar/\Delta t$ 数量级的精确度，其中的 $\Delta t$ 为两次测量之间的时间间隔.

(44.1)式通常称为能量的不确定度关系式.但是必须强调指出，它的含义与坐标和动量的不确定度关系式 $\Delta p\,\Delta x\sim\hbar$ 完全不同.在后一个关系式中，$\Delta p$ 和 $\Delta x$ 为同一时刻的动量和坐标的不确定度；它表明了这两个量不能同时具有定值.而另一方面，能量 $E,\varepsilon$ 却不论在哪一时刻都可以任意精确地进行测量.(44.1)式中的 $\left(E+\varepsilon\right)-\left(E'+\varepsilon'\right)$ 为能量 $E+\varepsilon$ 在两个不同时刻的精确测量值之差，并不是同一给定时刻的能量不确定度.

如果把 $E$ 看作某一系统的能量，把 $\varepsilon$ 看作“测量仪器”的能量，我们就可以说，它们之间的相互作用能量只能估计到 $\hbar/\Delta t$ 的程度为止.令 $\Delta E,\Delta\varepsilon,\cdots$ 为这些量的测量误差.在最好的情形下，当 $\varepsilon$ 和 $\varepsilon'$ 已精确知道时（$\Delta\varepsilon=\Delta\varepsilon'=0$），我们有
\begin{equation}
  \Delta\left(E-E'\right)\sim\frac{\hbar}{\Delta t}.
\end{equation}
根据这个式子，我们可以导出有关动量测量的若干重要结论.粒子（为确定起见，我们把它说成为电子）动量的测量过程，可以用该电子与某一其它粒子（称为“测量”粒子）的碰撞来实现，这个“测量”粒子在碰撞前后的动量假定可以精确地测定\footnote{在这里的分析中，这个“测量”粒子的能量等于多少并不重要.}.如果把动量守恒定律应用到这个碰撞问题中，可得三个方程式（一个矢量方程的三个分量式），具有六个未知数（碰撞前后电子动量的分量值）.如果该电子与许多“测量”粒子进行连续碰撞，并对每一次碰撞应用动量守恒定律，则方程数可以大大增加.可是未知数（碰撞中的电子动量）的数目也随之增加，并且易于证明，不管碰撞多少次，未知数的数目总比方程数多出三个.因此，为了测定电子的动量起见，除了动量守恒定律外，有必要在每一次碰撞中应用能量守恒定律.但是我们已经知道，应用后一个定律只能得到 $\hbar/\Delta t$ 数量级的精确度，$\Delta t$ 为测量过程始末的时间间隔.

为简化以后的讨论计，最好考虑一个假想的理想实验，其中的“测量粒子”是一个全反射的平面镜；这时重要的只是垂直于该平面镜的那个动量分量.为了确定粒子的动量 $P$，根据动量守恒定律和能量守恒定律给出下列方程：
\begin{equation}
  p'+P'-p-P=0,
\end{equation}
\begin{equation}
  \left|\varepsilon'+E'-\varepsilon-E\right|\sim\frac{\hbar}{\Delta t},
\end{equation}
式中的 $P,E$ 为粒子的动量和能量，$p,\varepsilon$ 为平面镜的动量和能量；不带撇的和带撇的分别标志碰撞前后时刻的量.“测量粒子”的 $p,p',\varepsilon,\varepsilon'$ 诸量，假定都是精确知道的，它们的误差都等于零.那么对其余各量的误差讲来，由上式得
\begin{equation*}
  \Delta P=\Delta P',\qquad
  \left|\Delta E'-\Delta E\right|\sim\frac{\hbar}{\Delta t}.
\end{equation*}
可是
\begin{equation*}
  \Delta E=\left(\partial E/\partial P\right)\Delta P=v\Delta P,
\end{equation*}
$v$ 为该电子（碰撞之前）的速度，同理有
\begin{equation*}
  \Delta E'=v'\Delta P'=v'\Delta P.
\end{equation*}
从而我们得
\begin{equation}
  \left|\left(v_x'-v_x\right)\Delta P_x\right|\sim\frac{\hbar}{\Delta t}.
\end{equation}
我们在速度和动量下面都加上了一个下标 $x$，这是为了强调这一式子对每一种分量分别讲来都是成立的.

这就是欲求的关系式.它表明电子动量的测量（具有给定的精确度 $\Delta P$）必然引起该电子的速度改变（亦即动量本身的改变）.测量过程的历时愈短，这种改变就愈大.只有 $\Delta t\to\infty$ 时，这种速度改变才能任意地小，但是占有长时间的动量测量只有对自由粒子才有意义.动量测量在经历短时间间隔后的不可重复性以及量子力学中测量的“两面性”——一个物理量的测量值与测量过程结束之后它所采取的数值之间加以区别的必要性——在这里表现得特别清楚\footnote{(44.5)式以及能量不确定度关系的物理含义是由 N.~玻尔给出的（1928）.}.

本节之初根据微扰论所得的结论，也可以用另一种观点推导出来，这就是考虑某一系统在某种微扰作用下的衰变问题.设 $E_0$ 为该系统的一个能级，这个能级是在完全略去衰变可能性的假定下计算出来的.令 $\tau$ 为系统处于这个态的寿命，也就是每单位时间衰变概率值的倒数.然后用同样的方法，我们可得
\begin{equation}
  \left|E_0-E-\varepsilon\right|\sim\frac{\hbar}{\tau},
\end{equation}
其中的 $E,\varepsilon$ 为该系统衰变之后形成的那两个部分的能量.$E+\varepsilon$ 可以看作衰变前的系统能量的一个估计值.因此所得的式子表明，一个处于“准稳定”态的可衰变系统的能量，只能确定到 $\hbar/\tau$ 数量级为止.这个量通常称为该能级的宽度 $\Gamma$.故
\begin{equation}
  \Gamma\sim\frac{\hbar}{\tau}.
\end{equation}

\section{以势能作微扰}

粒子在外场中的整个势能如果可以当作微扰来处理，这种情形值得专门考虑，此时的未受扰薛定谔方程就是该粒子的自由运动方程：
\begin{equation}
  \Delta\psi^{(0)}+k^2\psi^{(0)}=0,\qquad
  k=\frac{\sqrt{2mE}}{\hbar}=\frac{p}{\hbar},
\end{equation}
并具有平面波解.自由运动的能谱是连续的，因此我们现在所考虑的是连续谱微扰论中的一个独特情形.这个问题的直接求解，要比应用普遍公式更为方便.

一级近似的波函数修正 $\psi^{(1)}$ 满足下列方程：
\begin{equation}
  \Delta\psi^{(1)}+k^2\psi^{(1)}=\frac{2mU}{\hbar^2}\psi^{(0)},
\end{equation}
式中的 $U$ 为势能.根据电动力学我们知道，上式的解可以写成推迟势形式，即下列形式\footnote{这是(45.2)式的一个特解，其中还可以加进一个未受扰方程(45.1)式的任意解，即(45.2)式的右边等于零时的任意解.}：
\begin{equation}
  \psi^{(1)}(x,y,z)=-\frac{m}{2\pi\hbar^2}
  \int\psi^{(0)}U(x',y',z')\mathrm{e}^{ikr}\,\mathrm{d}V'/r,
\end{equation}
其中的
\begin{equation*}
  \mathrm{d}V'=\mathrm{d}x'\,\mathrm{d}y'\,\mathrm{d}z',\qquad
  r^2=\left(x-x'\right)^2+\left(y-y'\right)^2+\left(z-z'\right)^2.
\end{equation*}
现在来阐明势场 $U$ 应该满足什么样的条件才能看作微扰.微扰论的适用条件已经包含在 $\psi^{(1)}\ll\psi^{(0)}$ 这个要求内.设 $a$ 为该势场显著地不等于零的那个区域空间尺度的数量级.我们首先假定该粒子的能量十分小，以致 $ka$ 的数量级至多等于 1.此时(45.3)被积函数中的 $\mathrm{e}^{ikr}$ 因子在数量级的估计中并不重要，该积分的数量级为 $\psi^{(0)}|U|a^2$，故
\begin{equation*}
  \psi^{(1)}\sim\frac{ma^2|U|}{\hbar^2}\psi^{(0)}.
\end{equation*}
我们就得下列条件
\begin{equation}
  |U|\ll\frac{\hbar^2}{ma^2},\qquad\text{当}\ ka\lesssim1\ \text{时}.
\end{equation}
要注意的是，此式右边具有简单的物理含义；它等于封闭在线度为 $a$ 的体积内的一个粒子所具有的动能数量级（因为根据不确定度关系，粒子动量的数量级为 $\hbar/a$）.

我们来考虑一个特殊情形，设有一个十分浅的势阱，满足(45.4)式所给的条件.容易证明，在这样的势阱中并不存在负能级（R.~Peierls 1929）；对一个球对称的势阱讲来，这种特殊情形已经在 §33 的例题中证明过.实际上，当 $E=0$ 时，未受扰波函数变成一个常数，可以任意选定这个常数，我们暂定它等于 1；$\psi^{(0)}=1$.由于 $\psi^{(1)}\ll\psi^{(0)}$，在势阱中运动的波函数 $\psi=1+\psi^{(1)}$ 显然到处不等于零；这个本征函数是无节点的，它属于基态，因此 $E=0$ 仍是该粒子的最小能量值.由此可见，如果势阱足够浅，那么该粒子只能作无限运动；这个粒子不能被势阱所“俘获”.必须注意到，这是量子理论所特有的结果；经典力学中的粒子在任意势阱中都能作有限运动.

应该强调指出，以上所讲的只是针对三维势阱而言的.在一维或两维势阱中（即在这样的势阱中，$U$ 只是一个或两个坐标变量的函数），总可以存在负能级（见本节之末的例题）.原因在于，在一维或两维情形中，目前考虑的微扰论当能量 $E$ 等于零（或很小）时不再适用\footnote{在两维情形下，$\psi^{(1)}$ 可表为（根据二维波动方程的理论可知）类似于(45.3)的一个积分式，其中的 $\mathrm{e}^{ikr}\,\mathrm{d}x'\mathrm{d}y'\mathrm{d}z'/r$ 被换成 $i\pi\mathrm{H}_0^{(1)}(kr)\,\mathrm{d}x'\mathrm{d}y'$，$\mathrm{H}_0^{(1)}$ 为零阶的第一类汉克尔函数，而 $r^2=(x-x')^2+(y-y')^2$. 当 $k\to0$ 时，这个汉克尔函数，从而整个积分式，均以对数方式趋向无穷大.

同理，在一维情形下，被积函数中含有因子 $2\pi i\,\mathrm{e}^{ikr}\,\mathrm{d}x'/k$，其中的 $r=|x-x'|$，当 $k\to0$ 时，$\psi^{(1)}$ 按 $1/k$ 趋向无穷大.}.

在高能量情形下，此时有 $ka\gg1$，被积函数中的 $\mathrm{e}^{ikr}$ 因子起着重要作用，并使积分值显著地减小.这种情形下(45.3)式之解可以变换成另一种形式；为了推导这种形式，最好回到方程(45.2)直接求解.我们把微扰前的运动方向取作 $x$ 轴；则未受扰波函数的形式为 $\psi^{(0)}=\mathrm{e}^{ikx}$（常数因子暂定为 1）.我们来求下式之解：
\begin{equation*}
  \Delta\psi^{(1)}+k^2\psi^{(1)}=\frac{2m}{\hbar^2}U\mathrm{e}^{ikx},
\end{equation*}
令 $\psi^{(1)}=\mathrm{e}^{ikx}f$；鉴于所假定的 $k$ 值很大，我们在 $\Delta\psi^{(1)}$ 中只要保留 $\mathrm{e}^{ikx}$ 因子经过一次或两次微商后所得之项就可以了.由此可得 $f$ 的下列方程式：
\begin{equation*}
  2ik\,\frac{\partial f}{\partial x}=\frac{2mU}{\hbar^2},
\end{equation*}
故
\begin{equation}
  \psi^{(1)}=\mathrm{e}^{ikx}f=-\frac{im}{\hbar^2k}\,\mathrm{e}^{ikx}\int U\,\mathrm{d}x.
\end{equation}
对这个积分进行估计，得 $|\psi^{(1)}|\sim\dfrac{ma|U|}{\hbar^2k}$，因此这种情形下的微扰论适用条件为
\begin{equation}
  |U|\ll\frac{\hbar^2}{ma^2}ka=\frac{\hbar v}{a},\qquad ka\gg1,
\end{equation}
其中的 $v=k\hbar/m$ 是该粒子的速度.我们注意到这个条件比(45.4)来得弱.因此，在低能粒子情形下，如果一个势场可以当作微扰来处理，那么在高能情形下，这个势场一定仍能当作微扰来处理，反之就不一定正确了\footnote{一维情形下，微扰论的适用条件(45.6)式对所有的 $ka$ 成立. 三维情形下推得的(45.4)式对一维情形并不适用，因为所得的 $\psi^{(1)}$ 函数具有发散性（见前一附注）.}.

这里所讲的微扰论的适用性，对库仑场而言还需另行考虑.在 $U=\alpha/r$ 的场中，不可能分出一个有限空间区域，使得域外的 $U$ 比域内的小得很多.把(45.6)式中的参量 $a$ 改写成距离变量 $r$，就得欲求的适用条件；从而导出下列不等式：
\begin{equation}
  \frac{\alpha}{\hbar v}\ll1.
\end{equation}
由此可见，对高能粒子而言，库仑场可以当作微扰\footnote{必须注意，对 $U=\alpha/r$ 的场而言，(45.5)式当 $x/\sqrt{y^2+z^2}$ 很大时积分是发散的（对数发散）. 因此用微扰论求出的库仑场中的波函数，在绕 $x$ 轴的狭圆锥内不能适用.}.

最后，我们来推导一个公式，它能近似地给出某个粒子的波函数，只要这个粒子的能量 $E$ 到处都远大于势能 $U$（不加其它条件）.在一级近似中，这个波函数和坐标的关系，与自由运动（其方向取作 $x$ 轴）的情形相同.因此我们可设 $\psi$ 呈 $\psi=\mathrm{e}^{ikx}F$ 形式，其中的 $F$ 是一个坐标的函数，它与 $\mathrm{e}^{ikx}$ 相比变化得较慢（但是不能一般地认为它接近于 1）.把它代入薛定谔方程中，可得 $F$ 的方程式
\begin{equation}
  2ik\,\frac{\partial F}{\partial x}=\frac{2m}{\hbar^2}UF,
\end{equation}
由此得
\begin{equation}
  \psi=\mathrm{e}^{ikx}F=\text{常数}\times\mathrm{e}^{ikx}\,
  \mathrm{e}^{-\frac{i}{\hbar v}\int U\,\mathrm{d}x},
\end{equation}
这就是欲求的表式.但需注意，此式在远距离处并不适用.在方程(45.8)中，我们略去了含有 $F$ 二级微商的 $\Delta F$ 项.在远距离处，$\partial^2F/\partial x^2$ 和一级微商 $\partial F/\partial x$ 一起趋于零，但它对横坐标 $y,z$ 的偏微商并不趋于零，只有满足 $x\ll ka^2$ 条件时，它们才能被忽略掉.

\begin{xiti}

\noindent{\heiti 习题1}\quad 试求一维浅势阱中的能级.假定它已满足(45.4)式，并且积分 $\displaystyle\int_{-\infty}^{\infty}U\,\mathrm{d}x$ 是收敛的.

解：我们先假定能级 $|E|\ll|U|$，这个假定将被以后的结果所证实.那么，在下列薛定谔方程的右边
\begin{equation*}
  \frac{\mathrm{d}^2\psi}{\mathrm{d}x^2}=\frac{2m}{\hbar^2}\left[U(x)-E\right]\psi,
\end{equation*}
我们可以略去势阱范围内的 $E$，并把 $\psi$ 看作常数，不失其普遍性，可令 $\psi$ 等于 1：
\begin{equation*}
  \frac{\mathrm{d}^2\psi}{\mathrm{d}x^2}=\frac{2m}{\hbar^2}U.
\end{equation*}
此式对 $\mathrm{d}x$ 积分，积分限为 $\pm x_1$ 两点，并且 $a\ll x_1\ll1/\kappa$，其中的 $a$ 为势阱宽度，而 $\kappa=\sqrt{2m|E|}/\hbar$.由于 $U(x)$ 的积分是收敛的，式右的积分可延伸到从 $-\infty$ 到 $+\infty$ 的整个值域：
\setcounter{equation}{0}\renewcommand{\theequation}{\arabic{equation}}%
\begin{equation}
  \left.\frac{\mathrm{d}\psi}{\mathrm{d}x}\right|_{-x_1}^{x_1}
  =\frac{2m}{\hbar^2}\int_{-\infty}^{\infty}U\,\mathrm{d}x.
\end{equation}
在远离势阱处，波函数呈 $\psi=\mathrm{e}^{\pm\kappa x}$ 的形式.把它代入(1)式，我们得
\begin{equation*}
  -2\kappa=\frac{2m}{\hbar^2}\int_{-\infty}^{\infty}U\,\mathrm{d}x
\end{equation*}
或
\begin{equation*}
  |E|=\frac{m}{2\hbar^2}
  \left[\int_{-\infty}^{\infty}U\,\mathrm{d}x\right]^2.
\end{equation*}
由此可见，能级是一个很小的量，比阱深还高一个数量级（二级小量），与假设相符.

\noindent{\heiti 习题2}\quad 求二维浅势阱 $U=U(r)$ 中的能级（$r$ 为平面中的极坐标）；假定积分 $\displaystyle\int_0^{\infty}rU\,\mathrm{d}r$ 是收敛的.

解：和上题一样做法，我们得到阱内的方程为
\begin{equation*}
  \frac{1}{r}\frac{\mathrm{d}}{\mathrm{d}r}\left(r\frac{\mathrm{d}\psi}{\mathrm{d}r}\right)
  =\frac{2m}{\hbar^2}U.
\end{equation*}
此式对 $r$ 从 0 到 $r_1$ 积分（其中的 $a\ll r_1\ll1/\kappa$），结果得
\setcounter{equation}{0}%
\begin{equation}
  \left.\frac{\mathrm{d}\psi}{\mathrm{d}r}\right|_{r=r_1}
  =\frac{2m}{\hbar^2r_1}\int_0^{\infty}rU(r)\,\mathrm{d}r.
\end{equation}
在远离势阱处，两维自由运动的方程为
\begin{equation*}
  \frac{1}{r}\frac{\mathrm{d}}{\mathrm{d}r}\left(r\frac{\mathrm{d}\psi}{\mathrm{d}r}\right)
  +\frac{2m}{\hbar^2}E\psi=0,
\end{equation*}
它有（无穷远处为零的）解 $\psi=\text{常数}\times\mathrm{H}_0^{(1)}(i\kappa r)$；当这个函数的宗量值很小时，$\mathrm{H}_0^{(1)}$ 中的首项与 $\ln\kappa r$ 成正比.记住这一点，我们在 $r\sim a$ 处，把阱内外的 $\psi$ 的对数导数等同起来［阱内的 $\psi$ 对数导数等于(1)式的右边］，得
\begin{equation*}
  \frac{1}{a\ln\kappa a}\approx\frac{2m}{\hbar^2a}\int_0^{\infty}U(r)r\,\mathrm{d}r,
\end{equation*}
由此得
\begin{equation*}
  |E|\sim\frac{\hbar^2}{ma^2}\exp\biggl\{
  -\frac{\hbar^2}{m}\biggl|\int_0^{\infty}Ur\,\mathrm{d}r\biggr|^{-1}\biggr\}.
\end{equation*}
我们看到，能级比阱深（指数式地）小得很多.

\end{xiti}
\renewcommand{\theequation}{\arabic{section}.\arabic{equation}}%
